% ============================================================
%  MEMOIRE PFE - Plateforme securite routiere
%  Version complete corrigee
% ============================================================

\documentclass[a4paper,12pt]{report}

% --- Encodage et langue ---
\usepackage[utf8]{inputenc}
\usepackage[T1]{fontenc}
\usepackage[french]{babel}

% --- Mise en page ---
\usepackage[
  a4paper,
  top=3cm, bottom=2.5cm,
  left=2.5cm, right=2.5cm
]{geometry}

% --- Interligne ---
\usepackage{setspace}
\onehalfspacing

% --- Graphiques / TikZ ---
\usepackage{graphicx}
\usepackage{float}
\usepackage{tikz}
\usetikzlibrary{shapes.geometric, arrows.meta, positioning, fit, backgrounds, calc}

\graphicspath{{images/}{logos/}{diagrammes/}{captures/}{sites/}}

% --- Tableaux ---
\usepackage{booktabs}
\usepackage{tabularx}
\usepackage{array}
\usepackage{multirow}
\usepackage[table]{xcolor}
\renewcommand{\arraystretch}{1.45}

% --- Listings ---
\usepackage{listings}
\usepackage{inconsolata}

\definecolor{codeblue}{RGB}{0,0,200}
\definecolor{codegray}{RGB}{128,128,128}
\definecolor{codegreen}{RGB}{0,128,0}
\definecolor{backcolour}{RGB}{245,245,250}
\definecolor{codered}{RGB}{180,0,0}

\lstdefinestyle{mystyle}{
  backgroundcolor   = \color{backcolour},
  commentstyle      = \color{codegreen}\itshape,
  keywordstyle      = \color{codeblue}\bfseries,
  numberstyle       = \tiny\color{codegray},
  stringstyle       = \color{codered},
  basicstyle        = \ttfamily\footnotesize,
  breakatwhitespace = false,
  breaklines        = true,
  captionpos        = b,
  keepspaces        = true,
  numbers           = left,
  numbersep         = 8pt,
  showspaces        = false,
  showstringspaces  = false,
  showtabs          = false,
  tabsize           = 2,
  frame             = single,
  rulecolor         = \color{gray!40},
  xleftmargin       = 15pt,
  framexleftmargin  = 15pt
}
\lstset{style=mystyle}

% --- Boites ---
\usepackage{tcolorbox}
\tcbuselibrary{skins, breakable}
\newtcolorbox{notebox}{
  colback  = cyan!6,
  colframe = cyan!60!black,
  fonttitle= \bfseries,
  title    = Note,
  breakable
}

% --- En-tetes / pieds ---
\usepackage{fancyhdr}
\pagestyle{fancy}
\fancyhf{}
\fancyhead[R]{\small Chapitre \thechapter.\enspace\leftmark}
\fancyfoot[C]{\thepage}
\renewcommand{\headrulewidth}{0.4pt}
\renewcommand{\footrulewidth}{0pt}
\fancypagestyle{plain}{%
  \fancyhf{}%
  \fancyfoot[C]{\thepage}%
  \renewcommand{\headrulewidth}{0pt}%
}

% --- Titres ---
\usepackage{titlesec}
\titleformat{\chapter}[display]
  {\normalfont}
  {%
    \begin{tikzpicture}
      \draw[line width=0.5pt] (0,0) -- (1.2cm,0);
      \node[anchor=base west, inner sep=0pt] at (0, 0.15cm) {\footnotesize Chapitre};
      \node[anchor=base west, inner sep=0pt] at (1.4cm, -0.15cm)
            {\fontsize{60}{60}\selectfont\bfseries\thechapter};
      \draw[line width=0.5pt] (3.2cm, 0) -- (\textwidth, 0);
    \end{tikzpicture}%
  }
  {-10pt}
  {\huge\bfseries}
  []
\titlespacing*{\chapter}{0pt}{0pt}{30pt}

\titleformat{\section}
  {\normalfont\Large\bfseries}{\thesection}{1em}{}
\titlespacing*{\section}{0pt}{18pt}{8pt}

\titleformat{\subsection}
  {\normalfont\large\bfseries}{\thesubsection}{1em}{}
\titlespacing*{\subsection}{0pt}{14pt}{6pt}

\titleformat{\subsubsection}
  {\normalfont\normalsize\bfseries}{\thesubsubsection}{1em}{}
\titlespacing*{\subsubsection}{0pt}{10pt}{4pt}

% --- Minitoc ---
\usepackage{minitoc}
\setlength{\mtcindent}{0pt}
\mtcsetdepth{minitoc}{2}
\renewcommand{\mtctitle}{Sommaire}

% --- Captions ---
\usepackage{caption}
\captionsetup{
  labelfont = {sc, bf},
  labelsep  = endash,
  font      = small,
  textfont  = it
}

% --- Listes ---
\usepackage{enumitem}
\setlist[itemize]{label={---}, leftmargin=2em, itemsep=2pt, topsep=4pt}
\setlist[enumerate]{leftmargin=2em, itemsep=2pt, topsep=4pt}

% --- Hyperliens ---
\usepackage{hyperref}
\hypersetup{
  colorlinks = true,
  linkcolor  = red,
  urlcolor   = red,
  citecolor  = red,
  pdftitle   = {Memoire PFE -- Plateforme securite routiere},
  pdfauthor  = {Hammami Dhoha, Asma Ayari}
}

% --- Paragraphes ---
\setlength{\parindent}{1.5em}
\setlength{\parskip}{0pt}

% --- Macro sequence detaillee (avec zone erreurs optionnelle) ---
\newcommand{\ucseq}[8]{%
\begin{figure}[H]
  \centering
  \resizebox{0.96\textwidth}{!}{%
  \begin{tikzpicture}[x=1pt, y=1pt,
    box/.style={draw, fill=gray!10, minimum width=2.0cm, minimum height=0.65cm, font=\footnotesize, text centered},
    life/.style={draw, dashed, thin},
    msg/.style={draw, ->, >=Stealth, font=\scriptsize},
    ret/.style={draw, dashed, ->, >=Stealth, font=\scriptsize}
  ]
    \node[box] (A) at (0,0) {#2};
    \node[box] (F) at (120,0) {Frontend};
    \node[box] (API) at (250,0) {#3};
    \node[box] (S) at (390,0) {#4};

    \draw[life] (0,-10) -- (0,-300);
    \draw[life] (120,-10) -- (120,-300);
    \draw[life] (250,-10) -- (250,-300);
    \draw[life] (390,-10) -- (390,-300);

    \draw[msg] (0,-35) -- node[above]{#5} (120,-35);
    \draw[msg] (120,-65) -- node[above]{#6} (250,-65);
    \draw[msg] (250,-95) -- node[above]{traitement metier} (390,-95);

    % Fragment alt: branche erreur / branche nominale
    \draw[draw, dashed, rounded corners=3pt] (232,-112) rectangle (398,-268);
    \node[font=\scriptsize\bfseries, anchor=west] at (236,-118) {[alt]};

    % Branche erreur
    \draw[draw, dashed, rounded corners=2pt, fill=red!5] (238,-128) rectangle (392,-182);
    \node[font=\scriptsize\itshape, anchor=west] at (242,-135) {[erreur]};
    \node[font=\scriptsize, align=center] at (315,-160) {#7};
    \draw[ret] (250,-177) -- node[above]{4xx erreur} (120,-177);

    % Branche nominale
    \draw[draw, dashed, rounded corners=2pt, fill=green!5] (238,-190) rectangle (392,-258);
    \node[font=\scriptsize\itshape, anchor=west] at (242,-197) {[nominal]};
    \draw[ret] (390,-218) -- node[above]{resultat} (250,-218);
    \draw[ret] (250,-238) -- node[above]{#8} (120,-238);
    \draw[ret] (120,-262) -- node[above]{affichage UI} (0,-262);
  \end{tikzpicture}%
  }
  \caption{Diagramme de sequence -- #1}
\end{figure}
}

% ============================================================
\begin{document}
% ============================================================

\dominitoc
\setcounter{page}{2}

% ##########################################################
%  CHAPITRE 1 – Cadre general du projet
% ##########################################################
\chapter{Cadre general du projet}
\minitoc
\newpage

\section*{Introduction}
\addcontentsline{toc}{section}{Introduction}

Dans le cadre de notre stage de fin d'etudes au sein de l'Association Tunisie
pour la Securite Routiere (ATSR), nous avons ete amenes a concevoir et developper
une plateforme web intelligente dediee a la sensibilisation et a la formation
a la securite routiere.

Ce chapitre presente dans un premier temps l'organisme d'accueil, puis expose
la problematique identifiee, une etude comparative des solutions existantes,
la specification des besoins, la methodologie retenue, ainsi que l'architecture
et les technologies adoptees.

\section{Presentation de l'entreprise}

\begin{figure}[H]
  \centering
  \IfFileExists{logo_atsr.png}{%
    \includegraphics[width=0.4\textwidth]{logo_atsr.png}%
  }{%
    \begin{tikzpicture}
      \fill[gray!15] (0,0) rectangle (6cm,-3.2cm);
      \draw[gray!50] (0,0) rectangle (6cm,-3.2cm);
      \node[font=\small] at (3cm,-1.6cm) {Logo ATSR};
    \end{tikzpicture}%
  }
  \caption{Logo de l'Association Tunisie pour la Securite Routiere}
  \label{fig:logo_atsr}
\end{figure}

Notre projet a ete realise au sein de l'Association Tunisie pour la Securite
Routiere (ATSR). L'ATSR est une organisation a but non lucratif dont la vocation
est la sensibilisation, la formation et la promotion des bonnes pratiques en
matiere de securite routiere sur l'ensemble du territoire tunisien.

\begin{table}[H]
  \centering
  \caption{Fiche d'identite de l'organisme d'accueil}
  \label{tab:organisme}
  \begin{tabularx}{\textwidth}{>{\bfseries}p{4.5cm} X}
    \toprule
    Champ & Information \\
    \midrule
    Nom & Association Tunisie pour la Securite Routiere (ATSR) \\
    Type & Association a but non lucratif \\
    Adresse & Rue Ahmed Skolli -- Immeuble Alambra 2 -- 1\textsuperscript{er} etage -- 3027 Sfax \\
    Contact e-mail & \texttt{tsecuriteroutiere.sfax@gmail.com} \\
    Telephone & 20 412 543 \\
    Encadrant organisme & Lotfi Ghorbel \\
    Encadrant ISET & Elleuch Mohamed \\
    \bottomrule
  \end{tabularx}
\end{table}

    \subsection{Secteurs d'activite}

    L'association intervient dans plusieurs domaines complementaires lies a la
    securite routiere :

    \begin{itemize}
      \item Sensibilisation du grand public aux risques routiers via des campagnes
        de communication (affichages, medias, evenements).
      \item Formations et programmes pedagogiques a destination des auto-ecoles
        et des etablissements scolaires.
      \item Production et diffusion de contenus educatifs (videos, brochures,
        quiz et cours en ligne).
      \item Collaboration avec les autorites tunisiennes pour promouvoir le respect
        du code de la route.
      \item Collecte et analyse de donnees sur les accidents afin de cibler
        les actions de prevention.
    \end{itemize}

    \subsection{Organigramme de l'organisme d'accueil}

    \begin{figure}[H]
      \centering
      \begin{tikzpicture}
        \node[draw, fill=cyan!40, minimum width=3.2cm, minimum height=0.75cm] (pres) at (0,0) {President};
        \node[draw, fill=cyan!40, minimum width=3.2cm, minimum height=0.75cm] (dir)  at (0,-1.8) {Directeur Executif};
        \node[draw, fill=cyan!40, minimum width=3.2cm, minimum height=0.75cm] (cf)   at (-4.5,-3.6) {Commission Formation};
        \node[draw, fill=cyan!40, minimum width=3.2cm, minimum height=0.75cm] (cc)   at (0,-3.6) {Commission Communication};
        \node[draw, fill=cyan!40, minimum width=3.2cm, minimum height=0.75cm] (cp)   at (4.5,-3.6) {Commission Partenariats};
        \draw (pres) -- (dir);
        \draw (dir) -- (cf);
        \draw (dir) -- (cc);
        \draw (dir) -- (cp);
      \end{tikzpicture}
      \caption{Organigramme de l'organisme d'accueil}
      \label{fig:organigramme}
    \end{figure}

\section{Presentation du projet et etude de l'existant}

\subsection{Problematique}

La securite routiere represente un enjeu majeur pour la societe tunisienne.
Le nombre eleve d'accidents de la route, souvent cause par le non-respect du
code de la route ou par un manque de sensibilisation, justifie la mise en place
d'outils pedagogiques innovants et accessibles.

Les methodes traditionnelles d'apprentissage --- cours magistraux, brochures,
sessions en auto-ecole --- presentent plusieurs limites : faible interactivite,
absence de personnalisation, et difficulte d'acces a une assistance continue.

La problematique principale est donc :
\begin{quote}
\textit{Comment concevoir une plateforme intelligente et interactive permettant
d'enseigner la securite routiere, d'evaluer le niveau des apprenants de maniere
personnalisee et d'offrir une assistance pedagogique continue adaptee au contexte
tunisien ?}
\end{quote}

\subsection{Etude de l'existant}

Afin de positionner notre solution, nous avons etudie trois plateformes connues
du domaine : Permis-de-conduire.fr, iPERMIS et Moodle.

\subsubsection{Permis-de-conduire.fr}

Permis-de-conduire.fr est une plateforme web francaise dediee a la preparation
de l'examen theorique. Elle propose des cours en ligne classes par themes et des
QCM de simulation.

\begin{figure}[H]
  \centering
  \IfFileExists{sites/permis_fr.png}{%
    \includegraphics[width=0.86\textwidth]{sites/permis_fr.png}%
  }{%
    \begin{tikzpicture}
      \fill[gray!12] (0,0) rectangle (13cm,-5.5cm);
      \draw[gray!40] (0,0) rectangle (13cm,-5.5cm);
      \draw[gray!30] (0,0) -- (13cm,-5.5cm);
      \draw[gray!30] (13cm,0) -- (0,-5.5cm);
      \node[font=\small\itshape] at (6.5cm,-2.75cm) {[Capture Permis-de-conduire.fr]};
    \end{tikzpicture}%
  }
  \caption{Interface de Permis-de-conduire.fr}
  \label{fig:site_permis}
\end{figure}

\begin{table}[H]
  \centering
  \caption{Forces et faiblesses de Permis-de-conduire.fr}
  \label{tab:permis}
  \begin{tabularx}{\textwidth}{X X}
    \toprule
    \textbf{Forces} & \textbf{Faiblesses} \\
    \midrule
    Cours en ligne structures par themes.
      & Manque d'assistance intelligente integree. \\
    QCM de simulation d'examen.
      & Pas de support bilingue francais/arabe. \\
    Interface claire, disponible 24h/24.
      & Contenu non adapte au code de la route tunisien. \\
    \bottomrule
  \end{tabularx}
\end{table}

\subsubsection{iPERMIS}

iPERMIS est une application mobile de preparation du code de la route,
orientee vers la repetition des tests et le suivi de progression.

\begin{figure}[H]
  \centering
  \IfFileExists{sites/ipermis.png}{%
    \includegraphics[width=0.86\textwidth]{sites/ipermis.png}%
  }{%
    \begin{tikzpicture}
      \fill[gray!12] (0,0) rectangle (13cm,-5.5cm);
      \draw[gray!40] (0,0) rectangle (13cm,-5.5cm);
      \draw[gray!30] (0,0) -- (13cm,-5.5cm);
      \draw[gray!30] (13cm,0) -- (0,-5.5cm);
      \node[font=\small\itshape] at (6.5cm,-2.75cm) {[Capture iPERMIS]};
    \end{tikzpicture}%
  }
  \caption{Interface de l'application iPERMIS}
  \label{fig:site_ipermis}
\end{figure}

\begin{table}[H]
  \centering
  \caption{Forces et faiblesses d'iPERMIS}
  \label{tab:ipermis}
  \begin{tabularx}{\textwidth}{X X}
    \toprule
    \textbf{Forces} & \textbf{Faiblesses} \\
    \midrule
    Application mobile ergonomique.
      & Contenu uniquement en francais. \\
    Suivi de progression et memorisation par repetition.
      & Manque d'assistance intelligente. \\
    Disponible sur iOS et Android.
      & Pas de gestion documentaire pedagogique evoluee. \\
    \bottomrule
  \end{tabularx}
\end{table}

\subsubsection{Moodle}

Moodle est une plateforme LMS open source tres complete pour la gestion des
cours, quiz et activites d'apprentissage en ligne.

\begin{figure}[H]
  \centering
  \IfFileExists{sites/moodle.png}{%
    \includegraphics[width=0.86\textwidth]{sites/moodle.png}%
  }{%
    \begin{tikzpicture}
      \fill[gray!12] (0,0) rectangle (13cm,-5.5cm);
      \draw[gray!40] (0,0) rectangle (13cm,-5.5cm);
      \draw[gray!30] (0,0) -- (13cm,-5.5cm);
      \draw[gray!30] (13cm,0) -- (0,-5.5cm);
      \node[font=\small\itshape] at (6.5cm,-2.75cm) {[Capture Moodle]};
    \end{tikzpicture}%
  }
  \caption{Interface de Moodle}
  \label{fig:site_moodle}
\end{figure}

\begin{table}[H]
  \centering
  \caption{Forces et faiblesses de Moodle}
  \label{tab:moodle}
  \begin{tabularx}{\textwidth}{X X}
    \toprule
    \textbf{Forces} & \textbf{Faiblesses} \\
    \midrule
    Plateforme LMS complete et open source.
      & Manque d'assistance intelligente integree. \\
    Gestion des roles et quiz configurables.
      & Complexite de configuration elevee. \\
    Large communaute et ecosysteme de plugins.
      & Interface peu adaptee au grand public. \\
    \bottomrule
  \end{tabularx}
\end{table}

\subsubsection{Tableau comparatif des faiblesses}

\begin{table}[H]
  \centering
  \caption{Comparaison des limites des plateformes existantes}
  \label{tab:comparatif}
  \begin{tabularx}{\textwidth}{
    >{\raggedright\arraybackslash}X
    >{\centering\arraybackslash}p{2.8cm}
    >{\centering\arraybackslash}p{2.2cm}
    >{\centering\arraybackslash}p{2cm}
  }
    \toprule
    \textbf{Critere / Manque identifie}
      & \textbf{Permis-de-conduire.fr}
      & \textbf{iPERMIS}
      & \textbf{Moodle} \\
    \midrule
    Manque d'assistance intelligente
      & Oui & Oui & Oui \\
    Absence de support bilingue (FR/AR)
      & Oui & Oui & Oui \\
    Adaptation limitee au contexte tunisien
      & Oui & Oui & Partiel \\
    Module de detection visuelle des panneaux absent
      & Oui & Oui & Oui \\
    Interface difficile pour utilisateurs non techniques
      & Non & Non & Oui \\
    \bottomrule
  \end{tabularx}
\end{table}

\subsection{Solution proposee}

Notre solution consiste a developper une plateforme web intelligente dediee
a la securite routiere en Tunisie, avec les fonctionnalites suivantes :

\begin{itemize}
  \item Cours en ligne structures par thematiques.
  \item QCM de positionnement et d'evaluation avec correction automatique.
  \item Assistance conversationnelle bilingue (francais/arabe).
  \item Detection de panneaux de signalisation a partir d'images.
  \item Gestion des utilisateurs et des roles (invite, apprenant, formateur, administrateur).
  \item Tableaux de bord statistiques pour le suivi de l'activite.
  \item Architecture frontend/backend securisee et evolutive.
\end{itemize}

\section{Specification des besoins}

\subsection{Identification des acteurs}

\begin{enumerate}
  \item \textbf{Invite} : utilisateur non authentifie qui peut s'inscrire,
        se connecter (localement ou via Google) et lancer la recuperation de
        mot de passe, sans acces aux cours ni aux fonctionnalites protegees.
  \item \textbf{Apprenant} : utilisateur inscrit qui suit les cours,
        passe les QCM et interagit avec le chatbot.
  \item \textbf{Formateur} : responsable pedagogique qui cree et gere
        les cours.
  \item \textbf{Administrateur} : responsable fonctionnel et technique
        qui gere les utilisateurs, documents, QCM et statistiques.
  \item \textbf{Services externes} : OAuth Google, SMTP, services IA.
\end{enumerate}

\subsection{Besoins fonctionnels}

Les besoins fonctionnels decrivent les services attendus par chaque type
d'utilisateur. Ils constituent la base de priorisation du backlog et orientent
directement la conception des cas d'utilisation globaux et raffines.

\begin{table}[H]
  \centering
  \caption{Besoins fonctionnels par acteur}
  \label{tab:bf}
  \begin{tabular}{|p{3.2cm}|p{10cm}|}
    \hline
    \textbf{Acteur} & \textbf{Fonctionnalites} \\
    \hline
    Administrateur & Gestion des utilisateurs, documents, categories et QCM; statistiques; supervision chatbot. \\
    \hline
    Formateur & Creation/modification/publication de cours; upload videos; consultation des performances. \\
    \hline
    Apprenant & Consultation des cours; passation QCM; consultation resultats; chatbot; detection panneaux; gestion conversations. \\
    \hline
    Invite & Inscription; connexion locale/Google; mot de passe oublie; reinitialisation. \\
    \hline
  \end{tabular}
\end{table}

\subsection{Besoins non fonctionnels}

Les besoins non fonctionnels traduisent les exigences de qualite du systeme.
Ils garantissent qu'au-dela des fonctionnalites metier, la plateforme reste
fiable, securisee, performante et facilement maintenable dans la duree.

\begin{itemize}
  \item Performance
  \item Securite
  \item Disponibilite
  \item Scalabilite
  \item Maintenabilite
  \item Utilisabilite
  \item Internationalisation (FR/AR)
  \item Fiabilite
\end{itemize}

\subsection{Diagramme de cas d'utilisation global}

La lisibilite est amelioree en divisant le diagramme global sur deux pages.

% ------------------ Partie 1 ------------------
\begin{figure}[p]
\centering
\resizebox{0.96\textwidth}{!}{%
\begin{tikzpicture}[
  >=Stealth,
  actor/.style={draw, ellipse, minimum width=2.4cm, minimum height=0.9cm, font=\small},
  uc/.style={draw, ellipse, minimum width=4.8cm, minimum height=0.95cm, align=center, font=\small},
  rel/.style={draw, ->, line width=0.6pt}
]
  \draw[dashed, rounded corners=6pt] (0,0) rectangle (16,-11.8);
  \node[font=\small\bfseries] at (8,-0.5)
  {Plateforme Securite Routiere -- Cas d'utilisation globaux (Partie 1/2)};

  \node[actor] (G) at (1.8,-3.0) {Invite};
  \node[actor] (L) at (1.8,-8.2) {Apprenant};

  \node[uc] (UC1) at (8.2,-2.0) {S'inscrire};
  \node[uc] (UC2) at (8.2,-3.4) {Se connecter};
  \node[uc] (UC3) at (8.2,-4.8) {Consulter cours};
  \node[uc] (UC4) at (8.2,-6.2) {Passer un QCM};
  \node[uc] (UC5) at (8.2,-7.6) {Consulter resultats};
  \node[uc] (UC6) at (8.2,-9.0) {Discuter avec le chatbot};
  \node[uc] (UC7) at (8.2,-10.4) {Analyser panneau par image};

  % Invite: pas de consultation cours
  \draw[rel] (G) -- (UC1);
  \draw[rel] (G) -- (UC2);

  \draw[rel] (L) -- (UC2);
  \draw[rel] (L) -- (UC3);
  \draw[rel] (L) -- (UC4);
  \draw[rel] (L) -- (UC5);
  \draw[rel] (L) -- (UC6);
  \draw[rel] (L) -- (UC7);
\end{tikzpicture}%
}
\caption{Diagramme de cas d'utilisation global -- Partie 1 (Invite et Apprenant)}
\label{fig:uc_global_part1}
\end{figure}

\clearpage

% ------------------ Partie 2 ------------------
\begin{figure}[p]
\centering
\resizebox{0.96\textwidth}{!}{%
\begin{tikzpicture}[
  >=Stealth,
  actor/.style={draw, ellipse, minimum width=2.6cm, minimum height=0.9cm, font=\small},
  uc/.style={draw, ellipse, minimum width=5.0cm, minimum height=0.95cm, align=center, font=\small},
  rel/.style={draw, ->, line width=0.6pt}
]
  \draw[dashed, rounded corners=6pt] (0,0) rectangle (18,-12.8);
  \node[font=\small\bfseries] at (9,-0.5)
  {Plateforme Securite Routiere -- Cas d'utilisation globaux (Partie 2/2)};

  \node[actor] (F) at (1.8,-8.9) {Formateur};
  \node[actor] (A) at (1.8,-4.2) {Administrateur};

  \node[uc] (UC2)  at (8.0,-2.0)  {Se connecter};
  \node[uc] (UC3)  at (8.0,-3.4)  {Consulter cours};
  \node[uc] (UC8)  at (8.0,-4.8)  {Gerer utilisateurs};
  \node[uc] (UC9)  at (8.0,-6.2)  {Creer / modifier QCM};
  \node[uc] (UC10) at (8.0,-7.6)  {Publier QCM};
  \node[uc] (UC11) at (8.0,-9.0)  {Gerer categories QCM};
  \node[uc] (UC12) at (8.0,-10.4) {Gerer documents};
  \node[uc] (UC13) at (8.0,-11.8) {Lancer ingestion};

  \node[uc] (UC14) at (14.0,-4.8)  {Voir statistiques};
  \node[uc] (UC15) at (14.0,-7.0)  {Creer / modifier cours};
  \node[uc] (UC16) at (14.0,-8.8)  {Publier cours};
  \node[uc] (UC17) at (14.0,-10.6) {Uploader videos};
  \node[uc] (UC6)  at (14.0,-2.8)  {Discuter avec le chatbot};

  \draw[rel] (F) -- (UC2);
  \draw[rel] (F) -- (UC3);
  \draw[rel] (F) -- (UC15);
  \draw[rel] (F) -- (UC16);
  \draw[rel] (F) -- (UC17);

  \draw[rel] (A) -- (UC2);
  \draw[rel] (A) -- (UC3);
  \draw[rel] (A) -- (UC6);
  \draw[rel] (A) -- (UC8);
  \draw[rel] (A) -- (UC9);
  \draw[rel] (A) -- (UC10);
  \draw[rel] (A) -- (UC11);
  \draw[rel] (A) -- (UC12);
  \draw[rel] (A) -- (UC13);
  \draw[rel] (A) -- (UC14);
\end{tikzpicture}%
}
\caption{Diagramme de cas d'utilisation global -- Partie 2 (Formateur et Administrateur)}
\label{fig:uc_global_part2}
\end{figure}

\clearpage

\section{Methodologie}

\subsection{Planification du projet}

La planification a ete construite selon une approche iterative afin de livrer
des increments fonctionnels reguliers et validables avec les encadrants.
Le projet a ete decoupe en sprints courts pour reduire le risque, valider
rapidement les choix techniques et ajuster le perimetre en fonction des retours
metier et pedagogiques.

\subsection{Choix de l'approche Agile et de Scrum}

L'approche Agile a ete retenue car elle privilegie la collaboration continue,
la capacite d'adaptation et la livraison incrementale de valeur. Parmi les
methodes agiles, Scrum a ete selectionnee pour son cadre structure (roles,
ceremonies et artefacts) qui convient particulierement a un projet PFE mene sur
une duree limitee.

\begin{table}[H]
  \centering
  \caption{Comparaison des principales methodes agiles}
  \label{tab:methodes_agiles}
  \begin{tabularx}{\textwidth}{
    >{\bfseries\raggedright\arraybackslash}p{2.5cm}
    >{\raggedright\arraybackslash}X
    >{\raggedright\arraybackslash}X
    >{\raggedright\arraybackslash}X
  }
    \toprule
    Methode & Principe & Points forts & Limites \\
    \midrule
    Scrum
      & Sprints courts, backlog priorise, revues et retrospectives.
      & Cadre clair, livraisons frequentes, adaptation continue.
      & Necessite une coordination frequente de l'equipe. \\
    Kanban
      & Flux continu visualise sur tableau avec limites WIP.
      & Flexibilite elevee et bonne visibilite des taches.
      & Moins structure sur les jalons temporels. \\
    XP
      & TDD, pair programming, integration continue, refactoring.
      & Haute qualite logicielle et detection precoce des defauts.
      & Exige une forte maturite technique de l'equipe. \\
    \bottomrule
  \end{tabularx}
\end{table}

\textbf{Pourquoi Scrum :}
\begin{itemize}
  \item S'adapte bien a un projet PFE pilote en sprints.
  \item Facilite le suivi des priorites et la visibilite d'avancement.
  \item Permet l'integration rapide des retours des encadrants.
\end{itemize}

\subsection{Les roles dans Scrum}

Le cadre Scrum repose sur trois roles complementaires, chacun ayant une
responsabilite precise dans la reussite du projet :

\begin{itemize}
  \item \textbf{Product Owner} : priorise le backlog produit et garantit
  l'alignement des fonctionnalites avec les besoins metier.
  \item \textbf{Scrum Master} : facilite l'application de Scrum, supprime les
  blocages et accompagne l'equipe dans l'amelioration continue.
  \item \textbf{Equipe de developpement} : concoit, implemente, teste et livre
  l'increment fonctionnel a chaque sprint.
\end{itemize}

\subsection{Product Backlog}

Le Product Backlog regroupe l'ensemble des besoins exprimes sous forme de User
Stories. Chaque element est priorise selon sa valeur metier, son urgence et sa
faisabilite technique. Ce backlog reste evolutif et est raffine tout au long du
projet.

\begin{table}[H]
  \centering
  \caption{Backlog produit}
  \label{tab:backlog}
  \begin{tabular}{|c|p{8.2cm}|c|c|}
    \hline
    \textbf{ID} & \textbf{User Story} & \textbf{Priorite} & \textbf{Est.} \\
    \hline
    US1  & En tant que visiteur, je veux m'inscrire pour acceder a la plateforme & Haute & 3 SP \\
    \hline
    US2  & En tant qu'utilisateur, je veux me connecter avec Google & Moyenne & 5 SP \\
    \hline
    US3  & En tant qu'apprenant, je veux consulter les cours publies & Haute & 5 SP \\
    \hline
    US6  & En tant qu'admin, je veux creer et publier des QCM & Haute & 8 SP \\
    \hline
    US9  & En tant qu'utilisateur, je veux poser des questions a l'assistant & Haute & 8 SP \\
    \hline
    US12 & En tant qu'admin, je veux consulter les statistiques globales & Haute & 5 SP \\
    \hline
  \end{tabular}
\end{table}

\subsection{Planification des sprints}

La planification des sprints organise la mise en oeuvre progressive des modules
du systeme. Chaque sprint produit un increment testable afin de valider les
acquis avant de passer au sprint suivant.

\begin{table}[H]
  \centering
  \caption{Planification des sprints}
  \label{tab:sprints}
  \begin{tabular}{|c|p{4.5cm}|p{7cm}|}
    \hline
    \textbf{Sprint} & \textbf{Objectif} & \textbf{Livrables} \\
    \hline
    Sprint 1 & Base projet et authentification & Architecture initiale, auth locale, OAuth Google, roles \\
    \hline
    Sprint 2 & Module pedagogie & Gestion des cours, consultation apprenant \\
    \hline
    Sprint 3 & Evaluation & Gestion QCM, passation, scoring \\
    \hline
    Sprint 4 & Intelligence applicative & Assistant intelligent, ingestion, detection image \\
    \hline
    Sprint 5 & Administration et qualite & Dashboards, optimisation, tests \\
    \hline
    Sprint 6 & Finalisation & Deploiement, documentation, soutenance \\
    \hline
  \end{tabular}
\end{table}

\section{Architecture et outils de travail}

\subsection{Choix de l'architecture}

Le choix architectural devait concilier simplicite de mise en oeuvre,
maintenabilite et capacite d'evolution. Une architecture en couches a ete
privilegiee afin de separer clairement la presentation, la logique metier,
l'acces aux donnees et les composants IA.

\begin{table}[H]
  \centering
  \caption{Comparaison des architectures logicielles}
  \label{tab:archi_comp}
  \begin{tabularx}{\textwidth}{
    >{\bfseries\raggedright\arraybackslash}p{2.5cm}
    >{\raggedright\arraybackslash}X
    >{\raggedright\arraybackslash}X
    >{\raggedright\arraybackslash}X
  }
    \toprule
    Architecture & Principe & Avantages & Inconvenients \\
    \midrule
    Monolithique
      & Application unique qui regroupe toutes les fonctionnalites.
      & Mise en place initiale simple, deploiement unique.
      & Couplage fort et evolution difficile a grande echelle. \\
    En couches (Layered)
      & Separation presentation, logique metier, donnees et IA.
      & Maintenabilite, testabilite, responsabilites claires.
      & Demande une bonne discipline de conception. \\
    Microservices
      & Services independants deployables separement.
      & Scalabilite fine et isolation des composants.
      & Complexite operationnelle plus elevee. \\
    \bottomrule
  \end{tabularx}
\end{table}

\textbf{Choix retenu :} architecture en couches, adaptee a la taille du projet,
maintenable et evolutive.

\subsection{Architecture globale de la plateforme}

\begin{figure}[H]
  \centering
  \IfFileExists{architecture.png}{%
    \includegraphics[width=0.92\textwidth]{architecture.png}%
  }{%
    \begin{tikzpicture}
      \fill[gray!12] (0,0) rectangle (13.5cm,-6cm);
      \draw[gray!40] (0,0) rectangle (13.5cm,-6cm);
      \draw[gray!30] (0,0) -- (13.5cm,-6cm);
      \draw[gray!30] (13.5cm,0) -- (0,-6cm);
      \node[font=\small\itshape] at (6.75cm,-3cm) {[Schema architecture globale]};
    \end{tikzpicture}%
  }
  \caption{Architecture globale de la plateforme}
  \label{fig:archi_globale}
\end{figure}

\subsection{Architecture du RAG}

Le module de recherche augmentee suit cinq etapes : ingestion documentaire,
decoupage en segments, vectorisation, indexation, puis generation de reponse
contextualisee.

\begin{figure}[H]
  \centering
  \IfFileExists{architecture_rag.png}{%
    \includegraphics[width=0.92\textwidth]{architecture_rag.png}%
  }{%
    \begin{tikzpicture}
      \fill[gray!12] (0,0) rectangle (13.5cm,-6cm);
      \draw[gray!40] (0,0) rectangle (13.5cm,-6cm);
      \draw[gray!30] (0,0) -- (13.5cm,-6cm);
      \draw[gray!30] (13.5cm,0) -- (0,-6cm);
      \node[font=\small\itshape] at (6.75cm,-3cm) {[Schema architecture RAG]};
    \end{tikzpicture}%
  }
  \caption{Architecture du module de recherche augmentee}
  \label{fig:archi_rag}
\end{figure}

\subsection{Environnement de travail}

\begin{table}[H]
  \centering
  \caption{Environnement logiciel}
  \label{tab:env_logiciel}
  \begin{tabular}{|p{3cm}|p{4cm}|p{5cm}|}
    \hline
    \textbf{Outil} & \textbf{Technologie} & \textbf{Usage} \\
    \hline
    Systeme & Windows & Systeme d'exploitation \\
    \hline
    IDE & Visual Studio Code & Edition et debug \\
    \hline
    Frontend & Angular & Interface utilisateur \\
    \hline
    Backend & FastAPI (Python) & API REST \\
    \hline
    Base de donnees & MySQL & Stockage relationnel \\
    \hline
    Tests API & Postman & Validation endpoints \\
    \hline
    Versioning & GitHub & Gestion de versions \\
    \hline
  \end{tabular}
\end{table}

\subsection{Technologies utilisees}

\begin{table}[H]
  \centering
  \caption{Technologies principales du projet}
  \label{tab:techno}
  \begin{tabular}{|p{2.7cm}|p{2.2cm}|p{3.6cm}|p{4.3cm}|}
    \hline
    \textbf{Composant} & \textbf{Logo} & \textbf{Technologie} & \textbf{Justification} \\
    \hline
    Frontend & \IfFileExists{logos/angular.png}{\includegraphics[width=0.65cm]{logos/angular.png}}{Angular} & Angular, TypeScript & SPA modulaire, maintenable et responsive \\
    \hline
    Backend & \IfFileExists{logos/fastapi.png}{\includegraphics[width=0.65cm]{logos/fastapi.png}}{FastAPI} & FastAPI (Python) & API performante et integre facilement l'IA \\
    \hline
    Base de donnees & \IfFileExists{logos/mysql.png}{\includegraphics[width=0.7cm]{logos/mysql.png}}{MySQL} & MySQL + SQLAlchemy & Stockage relationnel fiable \\
    \hline
    Assistant intelligent & \IfFileExists{logos/llm.png}{\includegraphics[width=0.65cm]{logos/llm.png}}{IA} & LLM + recherche semantique & Assistance contextuelle aux apprenants \\
    \hline
    Authentification & \IfFileExists{logos/jwt.png}{\includegraphics[width=0.7cm]{logos/jwt.png}}{JWT} & JWT + OAuth Google + bcrypt & Securite et flexibilite d'acces \\
    \hline
    Outils & \IfFileExists{logos/github.png}{\includegraphics[width=0.65cm]{logos/github.png}}{GitHub} & GitHub, Postman, VS Code & Collaboration, tests, productivite \\
    \hline
  \end{tabular}
\end{table}

\subsection{Langages de programmation}

\begin{table}[H]
  \centering
  \caption{Langages de programmation utilises}
  \label{tab:langages}
  \begin{tabular}{|p{2.8cm}|p{2.2cm}|p{8cm}|}
    \hline
    \textbf{Langage} & \textbf{Logo} & \textbf{Usage dans le projet} \\
    \hline
    TypeScript & \IfFileExists{logos/typescript.png}{\includegraphics[width=0.65cm]{logos/typescript.png}}{TS} & Developpement du frontend Angular. \\
    \hline
    Python & \IfFileExists{logos/python.png}{\includegraphics[width=0.65cm]{logos/python.png}}{Py} & Developpement backend et logique metier. \\
    \hline
    SQL & \IfFileExists{logos/sql.png}{\includegraphics[width=0.65cm]{logos/sql.png}}{SQL} & Modelisation et requetes sur la base relationnelle. \\
    \hline
    HTML/CSS & \IfFileExists{logos/html_css.png}{\includegraphics[width=0.75cm]{logos/html_css.png}}{HTML/CSS} & Structuration et mise en forme des interfaces. \\
    \hline
  \end{tabular}
\end{table}

\subsection{Langage de modelisation}

Le langage UML est utilise pour modeliser le systeme a travers :

\begin{itemize}
  \item Diagrammes de cas d'utilisation pour les besoins fonctionnels.
  \item Diagrammes de classes pour la structure statique.
  \item Diagrammes de sequence pour les scenarios critiques.
\end{itemize}

\subsection{Choix technologiques detailles}

\begin{table}[H]
  \centering
  \caption{Choix technologiques detailles}
  \label{tab:choix_tech_detail}
  \begin{tabularx}{\textwidth}{>{\bfseries\raggedright\arraybackslash}p{2.8cm} p{3.8cm} X}
    \toprule
    Composant & Technologie & Justification \\
    \midrule
    Frontend & Angular (TypeScript) & Application SPA modulaire, maintenable et reactive. \\
    Backend & FastAPI (Python) & API REST performante et adaptee a l'integration IA. \\
    Base de donnees & MySQL + SQLAlchemy & Stockage relationnel robuste pour les donnees metier. \\
    Authentification & JWT + OAuth Google + bcrypt & Securisation des acces et gestion des roles. \\
    Assistant & LLM + recherche semantique & Assistance contextuelle basee sur la documentation. \\
    \bottomrule
  \end{tabularx}
\end{table}

\section*{Conclusion}
\addcontentsline{toc}{section}{Conclusion}

Ce chapitre a presente le cadre general, la problematique, les besoins,
la methodologie Scrum et l'architecture retenue pour la plateforme.


% ##########################################################
%  CHAPITRE 2 – Sprint 1
% ##########################################################
\chapter{Sprint 1 : Base projet et authentification}
\minitoc
\newpage

\section*{Introduction}
\addcontentsline{toc}{section}{Introduction}

Le Sprint 1 a pour objectif de poser les fondations techniques et fonctionnelles
de la plateforme. Il couvre :

\begin{enumerate}
  \item La mise en place de la base projet (backend FastAPI, frontend Angular,
  structure modulaire).
  \item Le socle de securite et d'authentification (JWT, OAuth Google,
  protection des acces).
  \item Les fonctions essentielles de gestion des comptes et des utilisateurs.
\end{enumerate}

Ce sprint constitue la base sur laquelle les sprints suivants (cours, QCM,
chatbot, statistiques) peuvent s'appuyer.

\section{Objectif et perimetre}

\subsection{Objectif general}

Fournir un systeme d'authentification fiable, securise et extensible, avec
gestion des roles (admin, formateur, apprenant).

\subsection{Perimetre fonctionnel retenu}

Les cas d'utilisation raffines pour ce sprint sont :

\begin{enumerate}
  \item \textbf{Authentification}
  \begin{itemize}
    \item Inscription locale
    \item Connexion locale
    \item Connexion Google
    \item Mot de passe oublie
    \item Reinitialisation du mot de passe
    \item Rafraichissement du token
  \end{itemize}
  \item \textbf{Gestion du profil}
  \begin{itemize}
    \item Consultation du profil
    \item Changement du mot de passe
  \end{itemize}
  \item \textbf{Gestion des utilisateurs (admin)}
  \begin{itemize}
    \item Lister les utilisateurs
    \item Creer un utilisateur
    \item Modifier le role
    \item Activer/desactiver un compte
    \item Supprimer un utilisateur
  \end{itemize}
\end{enumerate}

\section{Backlog Sprint 1}

\begin{table}[H]
  \centering
  \caption{Backlog du Sprint 1}
  \label{tab:backlog_s1}
  \begin{tabular}{|c|p{7.5cm}|c|c|}
    \hline
    \textbf{ID} & \textbf{User Story} & \textbf{Priorite} & \textbf{Est.} \\
    \hline
    US1 & En tant que visiteur, je veux m'inscrire & Haute & 3 SP \\
    \hline
    US2 & En tant que visiteur, je veux me connecter localement & Haute & 3 SP \\
    \hline
    US3 & En tant qu'utilisateur, je veux me connecter avec Google & Moyenne & 5 SP \\
    \hline
    US4 & En tant que visiteur, je veux reinitialiser mon mot de passe & Haute & 3 SP \\
    \hline
    US5 & En tant qu'utilisateur, je veux gerer mon profil & Moyenne & 3 SP \\
    \hline
    US6 & En tant qu'admin, je veux gerer les utilisateurs (CRUD, roles, statuts) & Haute & 8 SP \\
    \hline
  \end{tabular}
\end{table}

\section{Analyse}

\subsection{Acteurs du Sprint 1}
\begin{enumerate}
  \item Visiteur
  \item Utilisateur authentifie
  \item Administrateur
  \item Google OAuth
  \item Service Email SMTP
\end{enumerate}

\subsection{Diagramme de cas d'utilisation -- Sprint 1}

\begin{figure}[H]
  \centering
  \resizebox{0.94\textwidth}{!}{%
  \begin{tikzpicture}[
    >=Stealth,
    uc/.style={draw, ellipse, minimum width=4.0cm, minimum height=0.95cm, align=center, font=\footnotesize},
    actor/.style={font=\small\bfseries, align=center},
    rel/.style={draw, ->, line width=0.6pt},
    ext/.style={draw, dashed, ->, line width=0.6pt, font=\scriptsize\itshape}
  ]
    \draw[dashed, rounded corners=6pt] (0,0) rectangle (18,-15.2);
    \node[font=\small\bfseries] at (9,-0.5)
      {Systeme -- Sprint 1 : Base projet et authentification};

    \node[uc] (UC1)  at (9,-1.7)  {S'inscrire};
    \node[uc] (UC2)  at (9,-3.0)  {Se connecter\\localement};
    \node[uc] (UC3)  at (9,-4.3)  {Se connecter\\avec Google};
    \node[uc] (UC4)  at (9,-5.6)  {Mot de passe\\oublie};
    \node[uc] (UC5)  at (9,-6.9)  {Reinitialiser\\mot de passe};
    \node[uc] (UC6)  at (9,-8.2)  {Gerer profil\\(consulter / modifier)};
    \node[uc] (UC7)  at (9,-9.5)  {Changer\\mot de passe};
    \node[uc] (UC8)  at (9,-10.8) {Lister utilisateurs};
    \node[uc] (UC9)  at (9,-12.1) {Creer utilisateur};
    \node[uc] (UC10) at (9,-13.4) {Modifier role /\\activer-desactiver / suppression};

    \node[actor] (Vis) at (-1.8,-3.9)  {Visiteur};
    \node[actor] (Usr) at (-1.8,-8.9)  {Utilisateur\\authentifie};
    \node[actor] (Adm) at (-1.8,-12.1) {Administrateur};
    \node[actor] (Goo) at (19.8,-4.3)  {Google\\OAuth};
    \node[actor] (Eml) at (19.8,-5.6)  {Service\\Email SMTP};

    \draw[rel] (Vis) -- (UC1);
    \draw[rel] (Vis) -- (UC2);
    \draw[rel] (Vis) -- (UC3);
    \draw[rel] (Vis) -- (UC4);
    \draw[rel] (Vis) -- (UC5);

    \draw[rel] (Usr) -- (UC6);
    \draw[rel] (Usr) -- (UC7);

    \draw[rel] (Adm) -- (UC8);
    \draw[rel] (Adm) -- (UC9);
    \draw[rel] (Adm) -- (UC10);

    \draw[ext] (UC6) -- node[right]{$\ll$include$\gg$} (UC7);
    \draw[ext] (UC3) -- node[above]{verification token} (Goo);
    \draw[ext] (UC4) -- node[above]{envoi code} (Eml);
  \end{tikzpicture}%
  }
  \caption{Diagramme de cas d'utilisation -- Sprint 1 (corrige et centre)}
  \label{fig:uc_sprint1}
\end{figure}

\section{Conception}

\subsection{Cas d'utilisation raffines}

\subsubsection{UC1 -- S'inscrire}
\begin{table}[H]
  \centering
  \caption{Description du cas d'utilisation UC1}
  \begin{tabularx}{\textwidth}{>{\bfseries}p{4cm} X}
    \toprule
    Cas d'utilisation & S'inscrire localement \\
    \midrule
    Acteur principal & Visiteur \\
    Preconditions & Email et username non utilises. \\
    Scenario nominal & Validation des champs, verification unicite, hachage du mot de passe, creation compte, generation JWT. \\
    Postconditions & Compte actif cree et session ouverte. \\
    Exceptions & Email/username deja existants, mot de passe non conforme. \\
    \bottomrule
  \end{tabularx}
\end{table}

\ucseq{UC1 -- S'inscrire}{Visiteur}{API Auth}{AuthService + Base}{Saisie formulaire inscription}{POST /api/auth/register}{Email deja utilise ou mot de passe invalide (HTTP 400)}{201 Compte cree + JWT}

\subsubsection{UC2 -- Se connecter localement}
\begin{table}[H]
  \centering
  \caption{Description du cas d'utilisation UC2}
  \begin{tabularx}{\textwidth}{>{\bfseries}p{4cm} X}
    \toprule
    Cas d'utilisation & Se connecter localement \\
    \midrule
    Acteur principal & Visiteur \\
    Preconditions & Compte local existant et actif. \\
    Scenario nominal & Verification identifiants, verification hash mot de passe, generation des tokens. \\
    Postconditions & Session authentifiee ouverte. \\
    Exceptions & Identifiants invalides, compte desactive, blocage rate limit. \\
    \bottomrule
  \end{tabularx}
\end{table}

\ucseq{UC2 -- Se connecter localement}{Visiteur}{API Auth}{AuthService + Base}{Saisie email + mot de passe}{POST /api/auth/login}{Mot de passe incorrect / token absent ou invalide (401/403)}{200 TokenResponse}

\subsubsection{UC3 -- Se connecter avec Google}
\begin{table}[H]
  \centering
  \caption{Description du cas d'utilisation UC3}
  \begin{tabularx}{\textwidth}{>{\bfseries}p{4cm} X}
    \toprule
    Cas d'utilisation & Se connecter avec Google \\
    \midrule
    Acteur principal & Visiteur \\
    Acteur secondaire & Google OAuth \\
    Preconditions & Credential Google valide. \\
    Scenario nominal & Verification du token Google, creation/liaison compte local, generation JWT. \\
    Postconditions & Session ouverte. \\
    Exceptions & Token Google invalide, audience non autorisee. \\
    \bottomrule
  \end{tabularx}
\end{table}

\ucseq{UC3 -- Se connecter avec Google}{Visiteur}{API Auth}{Google OAuth}{Clic sur connexion Google}{POST /api/auth/google}{Token Google invalide ou email non verifie (401)}{200 Session ouverte + JWT}

\subsubsection{UC4 -- Mot de passe oublie}
\begin{table}[H]
  \centering
  \caption{Description du cas d'utilisation UC4}
  \begin{tabularx}{\textwidth}{>{\bfseries}p{4cm} X}
    \toprule
    Cas d'utilisation & Mot de passe oublie \\
    \midrule
    Acteur principal & Visiteur \\
    Acteur secondaire & Service Email SMTP \\
    Preconditions & Adresse email saisie. \\
    Scenario nominal & Generation d'un code temporaire et envoi SMTP avec message neutre. \\
    Postconditions & Code stocke avec expiration. \\
    Exceptions & Erreur SMTP ou email inexistant (message neutre conserve). \\
    \bottomrule
  \end{tabularx}
\end{table}

\ucseq{UC4 -- Mot de passe oublie}{Visiteur}{API Auth}{Email Service}{Saisie email}{POST /api/auth/forgot-password}{Erreur SMTP ou compte inexistant (message neutre)}{200 Message neutre}

\subsubsection{UC5 -- Reinitialiser le mot de passe}
\begin{table}[H]
  \centering
  \caption{Description du cas d'utilisation UC5}
  \begin{tabularx}{\textwidth}{>{\bfseries}p{4cm} X}
    \toprule
    Cas d'utilisation & Reinitialiser le mot de passe \\
    \midrule
    Acteur principal & Visiteur \\
    Preconditions & Code valide et non expire. \\
    Scenario nominal & Verification code, hachage du nouveau mot de passe, mise a jour et invalidation du code. \\
    Postconditions & Nouveau mot de passe actif. \\
    Exceptions & Code invalide/expire, utilisateur introuvable. \\
    \bottomrule
  \end{tabularx}
\end{table}

\ucseq{UC5 -- Reinitialiser le mot de passe}{Visiteur}{API Auth}{AuthService + Base}{Saisie email + code + nouveau MDP}{POST /api/auth/reset-password}{Code incorrect/expire (400)}{200 Password reset}

\subsubsection{UC6 -- Gerer son profil}
\begin{table}[H]
  \centering
  \caption{Description du cas d'utilisation UC6}
  \begin{tabularx}{\textwidth}{>{\bfseries}p{4cm} X}
    \toprule
    Cas d'utilisation & Gerer son profil (consulter / modifier) \\
    \midrule
    Acteur principal & Utilisateur authentifie \\
    Preconditions & Token JWT valide. \\
    Scenario nominal & Consultation des donnees, mise a jour des informations du profil. \\
    Postconditions & Profil affiche et/ou mis a jour. \\
    Exceptions & Token invalide/expire, compte desactive. \\
    \bottomrule
  \end{tabularx}
\end{table}

\ucseq{UC6 -- Gerer son profil}{Utilisateur authentifie}{API Auth}{Base MySQL}{Acces page profil}{GET/PUT /api/auth/me}{Token JWT invalide ou expire (401)}{200 Profil retourne/mis a jour}

\subsubsection{UC7 -- Changer son mot de passe}
\begin{table}[H]
  \centering
  \caption{Description du cas d'utilisation UC7}
  \begin{tabularx}{\textwidth}{>{\bfseries}p{4cm} X}
    \toprule
    Cas d'utilisation & Changer son mot de passe \\
    \midrule
    Acteur principal & Utilisateur authentifie \\
    Preconditions & Ancien mot de passe correct. \\
    Scenario nominal & Verification ancien MDP, hachage nouveau MDP et mise a jour. \\
    Postconditions & Nouveau mot de passe actif. \\
    Exceptions & Ancien MDP incorrect, nouveau MDP invalide, compte Google pur. \\
    \bottomrule
  \end{tabularx}
\end{table}

\ucseq{UC7 -- Changer son mot de passe}{Utilisateur authentifie}{API Auth}{AuthService + Base}{Saisie ancien + nouveau MDP}{POST /api/auth/change-password}{Ancien mot de passe incorrect (400)}{200 Password updated}

\subsubsection{UC8 -- Lister les utilisateurs}
\begin{table}[H]
  \centering
  \caption{Description du cas d'utilisation UC8}
  \begin{tabularx}{\textwidth}{>{\bfseries}p{4cm} X}
    \toprule
    Cas d'utilisation & Lister les utilisateurs \\
    \midrule
    Acteur principal & Administrateur \\
    Preconditions & Utilisateur connecte avec role admin. \\
    Scenario nominal & Consultation liste paginee avec filtres. \\
    Postconditions & Vue de supervision actualisee. \\
    Exceptions & Acces refuse si role non admin. \\
    \bottomrule
  \end{tabularx}
\end{table}

\ucseq{UC8 -- Lister les utilisateurs}{Administrateur}{API Users}{Base MySQL}{Demande liste paginee}{GET /api/users}{Token invalide ou role insuffisant (401/403)}{200 Liste utilisateurs}

\subsubsection{UC9 -- Creer un utilisateur}
\begin{table}[H]
  \centering
  \caption{Description du cas d'utilisation UC9}
  \begin{tabularx}{\textwidth}{>{\bfseries}p{4cm} X}
    \toprule
    Cas d'utilisation & Creer un utilisateur \\
    \midrule
    Acteur principal & Administrateur \\
    Preconditions & Role admin et donnees valides. \\
    Scenario nominal & Verification unicite, hachage mot de passe, creation compte. \\
    Postconditions & Nouvel utilisateur present dans le systeme. \\
    Exceptions & Email/username existant, role invalide. \\
    \bottomrule
  \end{tabularx}
\end{table}

\ucseq{UC9 -- Creer un utilisateur}{Administrateur}{API Users}{AuthService + Base}{Saisie donnees utilisateur}{POST /api/users}{Email/username deja existant (400)}{201 Utilisateur cree}

\subsubsection{UC10 -- Modifier role / activer-desactiver / supprimer}
\begin{table}[H]
  \centering
  \caption{Description du cas d'utilisation UC10}
  \begin{tabularx}{\textwidth}{>{\bfseries}p{4cm} X}
    \toprule
    Cas d'utilisation & Modifier role / activer-desactiver / suppression \\
    \midrule
    Acteur principal & Administrateur \\
    Preconditions & Role admin et utilisateur cible existant. \\
    Scenario nominal & Changement role, activation/desactivation ou suppression selon regles metier. \\
    Postconditions & Compte cible mis a jour. \\
    Exceptions & Auto-modification critique interdite, utilisateur non trouve. \\
    \bottomrule
  \end{tabularx}
\end{table}

\ucseq{UC10 -- Modifier role / activer-desactiver / suppression}{Administrateur}{API Users}{Base MySQL}{Choix utilisateur + action admin}{PUT/DELETE /api/users/...}{Auto-modification interdite ou cible absente}{200 Action appliquee}

\section{Diagrammes de sequence detailles du Sprint 1}

\subsection{Sequence detaillee -- Connexion locale}

\begin{figure}[H]
  \centering
  \begin{tikzpicture}[x=1pt, y=1pt,
    box/.style={draw, fill=gray!10, minimum width=1.8cm, minimum height=0.6cm, font=\footnotesize},
    msg/.style={draw, ->, >=Stealth, font=\scriptsize},
    ret/.style={draw, dashed, ->, >=Stealth, font=\scriptsize},
    life/.style={draw, dashed, thin}
  ]
  \foreach \x/\lbl in {0/Visiteur, 90/Frontend, 190/API Auth, 290/AuthService, 390/Base MySQL}{
    \node[box] at (\x,0) {\lbl};
    \draw[life] (\x,-8) -- (\x,-320);
  }
  \draw[msg] (0,-30) -- node[above]{email + mot de passe} (90,-30);
  \draw[msg] (90,-55) -- node[above]{POST /api/auth/login} (190,-55);
  \draw[msg] (190,-80) -- node[above]{charger utilisateur} (390,-80);
  \draw[ret] (390,-100) -- node[above]{User} (190,-100);
  \draw[msg] (190,-125) -- node[above]{verify\_password()} (290,-125);

  \draw[draw, dashed, rounded corners=3pt] (175,-145) rectangle (405,-240);
  \node[font=\scriptsize\bfseries] at (195,-152) {[alt]};
  \node[font=\scriptsize\itshape] at (235,-172) {[mot de passe invalide]};
  \draw[ret] (190,-190) -- node[above]{401 Invalid credentials} (90,-190);
  \node[font=\scriptsize\itshape] at (235,-212) {[mot de passe valide]};
  \draw[msg] (190,-228) -- node[above]{create tokens JWT} (290,-228);
  \draw[ret] (290,-248) -- node[above]{access + refresh} (190,-248);

  \draw[ret] (190,-275) -- node[above]{200 TokenResponse} (90,-275);
  \draw[ret] (90,-300) -- node[above]{session ouverte} (0,-300);
  \end{tikzpicture}
  \caption{Diagramme de sequence detaille -- Connexion locale}
  \label{fig:seq_login_detail}
\end{figure}

\subsection{Sequence detaillee -- Reinitialisation du mot de passe}

\begin{figure}[H]
  \centering
  \begin{tikzpicture}[x=1pt, y=1pt,
    box/.style={draw, fill=gray!10, minimum width=1.8cm, minimum height=0.6cm, font=\footnotesize},
    msg/.style={draw, ->, >=Stealth, font=\scriptsize},
    ret/.style={draw, dashed, ->, >=Stealth, font=\scriptsize},
    life/.style={draw, dashed, thin}
  ]
  \foreach \x/\lbl in {0/Visiteur, 95/Frontend, 200/API Auth, 305/AuthService, 410/Base MySQL}{
    \node[box] at (\x,0) {\lbl};
    \draw[life] (\x,-8) -- (\x,-340);
  }
  \draw[msg] (0,-30) -- node[above]{email + code + nouveau MDP} (95,-30);
  \draw[msg] (95,-55) -- node[above]{POST /api/auth/reset-password} (200,-55);
  \draw[msg] (200,-80) -- node[above]{verify\_reset\_code()} (305,-80);

  \draw[draw, dashed, rounded corners=3pt] (185,-100) rectangle (425,-235);
  \node[font=\scriptsize\bfseries] at (205,-107) {[alt]};
  \node[font=\scriptsize\itshape] at (250,-128) {[code invalide/expire]};
  \draw[ret] (200,-146) -- node[above]{400 Invalid or expired code} (95,-146);
  \node[font=\scriptsize\itshape] at (250,-170) {[code valide]};
  \draw[msg] (200,-188) -- node[above]{hash\_password(new\_password)} (305,-188);
  \draw[ret] (305,-208) -- node[above]{hashed\_password} (200,-208);
  \draw[msg] (200,-228) -- node[above]{update user + commit} (410,-228);

  \draw[ret] (410,-255) -- node[above]{OK} (200,-255);
  \draw[ret] (200,-282) -- node[above]{200 Password reset} (95,-282);
  \draw[ret] (95,-310) -- node[above]{confirmation} (0,-310);
  \end{tikzpicture}
  \caption{Diagramme de sequence detaille -- Reinitialisation du mot de passe}
  \label{fig:seq_reset_detail}
\end{figure}

\subsection{Sequence detaillee -- Administration des utilisateurs}

\begin{figure}[H]
  \centering
  \begin{tikzpicture}[x=1pt, y=1pt,
    box/.style={draw, fill=gray!10, minimum width=1.8cm, minimum height=0.6cm, font=\footnotesize},
    msg/.style={draw, ->, >=Stealth, font=\scriptsize},
    ret/.style={draw, dashed, ->, >=Stealth, font=\scriptsize},
    life/.style={draw, dashed, thin}
  ]
  \foreach \x/\lbl in {0/Admin, 95/Users UI, 200/API Users, 305/AuthService, 410/Base MySQL}{
    \node[box] at (\x,0) {\lbl};
    \draw[life] (\x,-8) -- (\x,-330);
  }
  \draw[msg] (0,-30) -- node[above]{choix user + action} (95,-30);
  \draw[msg] (95,-55) -- node[above]{PUT /api/users/{id}/role} (200,-55);
  \draw[msg] (200,-80) -- node[above]{verify token + role admin} (305,-80);
  \draw[ret] (305,-100) -- node[above]{OK admin} (200,-100);
  \draw[msg] (200,-125) -- node[above]{load target user} (410,-125);
  \draw[ret] (410,-145) -- node[above]{User cible} (200,-145);

  \draw[draw, dashed, rounded corners=3pt] (185,-165) rectangle (425,-255);
  \node[font=\scriptsize\bfseries] at (205,-172) {[alt]};
  \node[font=\scriptsize\itshape] at (255,-192) {[auto-modification critique]};
  \draw[ret] (200,-210) -- node[above]{400 Action interdite} (95,-210);
  \node[font=\scriptsize\itshape] at (250,-233) {[cas nominal]};
  \draw[msg] (200,-248) -- node[above]{update + commit} (410,-248);

  \draw[ret] (410,-273) -- node[above]{OK} (200,-273);
  \draw[ret] (200,-300) -- node[above]{200 user mis a jour} (95,-300);
  \end{tikzpicture}
  \caption{Diagramme de sequence detaille -- Administration des utilisateurs}
  \label{fig:seq_admin_detail}
\end{figure}

\subsection{Extraits techniques du Sprint 1}

\textbf{Exemple 1 -- Inscription avec hachage du mot de passe}
\begin{lstlisting}[language=Python, caption={POST /api/auth/register}]
@router.post("/register", response_model=TokenResponse)
async def register(req: RegisterRequest, db: Session = Depends(get_db)):
    if db.query(User).filter(User.email == req.email).first():
        raise HTTPException(status_code=400, detail="Email already registered")
    user = User(
        username=req.username,
        email=req.email,
        hashed_password=hash_password(req.password),
        role=UserRole.apprenant,
        is_active=True
    )
    db.add(user)
    db.commit()
    db.refresh(user)
    access = create_access_token({"sub": str(user.id)})
    refresh = create_refresh_token({"sub": str(user.id)})
    return TokenResponse(access_token=access, refresh_token=refresh, user=user)
\end{lstlisting}

\textbf{Exemple 2 -- Connexion locale et verification mot de passe}
\begin{lstlisting}[language=Python, caption={POST /api/auth/login}]
@router.post("/login", response_model=TokenResponse)
async def login(req: LoginRequest, db: Session = Depends(get_db)):
    user = db.query(User).filter(User.email == req.email).first()
    if not user or not verify_password(req.password, user.hashed_password):
        raise HTTPException(status_code=401, detail="Invalid credentials")
    access = create_access_token({"sub": str(user.id), "role": user.role})
    refresh = create_refresh_token({"sub": str(user.id)})
    return TokenResponse(access_token=access, refresh_token=refresh, user=user)
\end{lstlisting}

\textbf{Exemple 3 -- Reinitialisation du mot de passe}
\begin{lstlisting}[language=Python, caption={POST /api/auth/reset-password}]
@router.post("/reset-password")
async def reset_password(req: ResetPasswordRequest, db: Session = Depends(get_db)):
    entry = reset_codes.get(req.email)
    if not entry or entry["code"] != req.code:
        raise HTTPException(status_code=400, detail="Invalid code")
    if datetime.utcnow() > entry["expires"]:
        raise HTTPException(status_code=400, detail="Code expired")
    user = db.query(User).filter(User.email == req.email).first()
    user.hashed_password = hash_password(req.new_password)
    db.commit()
    del reset_codes[req.email]
    return {"message": "Password successfully reset"}
\end{lstlisting}

\subsection{Diagramme de classes -- Sprint 1}

\begin{figure}[H]
  \centering
  \begin{tikzpicture}[
    cls/.style={draw, rounded corners=2pt, fill=gray!10, align=left, font=\scriptsize, text width=4.6cm, inner sep=4pt},
    dep/.style={draw, dashed, ->, >=Stealth, line width=0.6pt},
    assoc/.style={draw, ->, >=Stealth, line width=0.6pt}
  ]

  \node[cls] (User) at (0,0) {
    \textbf{User}\\
    \rule{4.2cm}{0.3pt}\\
    id: int\\
    username: string\\
    email: string\\
    hashed\_password: string\\
    full\_name: string\\
    role: UserRole\\
    is\_active: bool
  };

  \node[cls] (Role) at (8.2,0) {
    \textbf{<<enumeration>> UserRole}\\
    \rule{4.2cm}{0.3pt}\\
    admin\\
    formateur\\
    apprenant
  };

  \node[cls] (AuthAPI) at (0,-4.8) {
    \textbf{AuthAPI}\\
    \rule{4.2cm}{0.3pt}\\
    register(req)\\
    login(req)\\
    google\_auth(req)\\
    forgot\_password(req)\\
    reset\_password(req)\\
    get\_profile()
  };

  \node[cls] (AuthService) at (8.2,-4.8) {
    \textbf{AuthService}\\
    \rule{4.2cm}{0.3pt}\\
    hash\_password(p)\\
    verify\_password(p,h)\\
    create\_access\_token(d)\\
    decode\_token(t)
  };

  \node[cls] (UsersAPI) at (4.1,-9.3) {
    \textbf{UsersAPI}\\
    \rule{4.2cm}{0.3pt}\\
    list\_users(...)\\
    create\_user(req)\\
    update\_user\_role(id)\\
    delete\_user(id)
  };

  \draw[assoc] (User.east) -- node[above]{role} (Role.west);
  \draw[dep] (AuthAPI.east) -- node[above]{utilise} (AuthService.west);
  \draw[dep] (UsersAPI.east) -- node[right]{utilise} (AuthService.south);
  \draw[dep] (AuthAPI.north) -- node[left]{persiste} (User.south);
  \draw[dep] (UsersAPI.west) -- node[left]{persiste} (User.south east);
  \end{tikzpicture}
\caption{Diagramme de classes -- Sprint 1}
\label{fig:class_sprint1}
\end{figure}

\subsection{Captures d'ecran des interfaces principales}

\begin{figure}[H]
  \centering
  \IfFileExists{captures/register.png}{%
    \includegraphics[width=0.85\textwidth]{captures/register.png}%
  }{%
    \begin{tikzpicture}
      \fill[gray!12] (0,0) rectangle (13cm,-5.8cm);
      \draw[gray!40] (0,0) rectangle (13cm,-5.8cm);
      \node[font=\small\itshape] at (6.5cm,-2.9cm) {[Capture page inscription]};
    \end{tikzpicture}%
  }
  \caption{Interface d'inscription}
\end{figure}

\begin{figure}[H]
  \centering
  \IfFileExists{captures/login.png}{%
    \includegraphics[width=0.85\textwidth]{captures/login.png}%
  }{%
    \begin{tikzpicture}
      \fill[gray!12] (0,0) rectangle (13cm,-5.8cm);
      \draw[gray!40] (0,0) rectangle (13cm,-5.8cm);
      \node[font=\small\itshape] at (6.5cm,-2.9cm) {[Capture page connexion]};
    \end{tikzpicture}%
  }
  \caption{Interface de connexion}
\end{figure}

\begin{figure}[H]
  \centering
  \IfFileExists{captures/profile.png}{%
    \includegraphics[width=0.85\textwidth]{captures/profile.png}%
  }{%
    \begin{tikzpicture}
      \fill[gray!12] (0,0) rectangle (13cm,-5.8cm);
      \draw[gray!40] (0,0) rectangle (13cm,-5.8cm);
      \node[font=\small\itshape] at (6.5cm,-2.9cm) {[Capture gestion profil]};
    \end{tikzpicture}%
  }
  \caption{Interface de gestion du profil}
\end{figure}

\begin{figure}[H]
  \centering
  \IfFileExists{captures/admin_users.png}{%
    \includegraphics[width=0.90\textwidth]{captures/admin_users.png}%
  }{%
    \begin{tikzpicture}
      \fill[gray!12] (0,0) rectangle (13.5cm,-6.2cm);
      \draw[gray!40] (0,0) rectangle (13.5cm,-6.2cm);
      \node[font=\small\itshape] at (6.75cm,-3.1cm) {[Capture administration utilisateurs]};
    \end{tikzpicture}%
  }
  \caption{Interface d'administration des utilisateurs}
\end{figure}

\section{Evaluation du Sprint 1}

\subsection{Scenarios de tests realises}

\begin{table}[H]
  \centering
  \caption{Table de tests fonctionnels du Sprint 1}
  \label{tab:tests_s1}
  \begin{tabular}{|p{5.2cm}|p{3.2cm}|p{5.4cm}|}
    \hline
    \textbf{Scenario} & \textbf{Resultat attendu} & \textbf{Statut} \\
    \hline
    Inscription avec donnees valides & Compte cree + JWT & Valide \\
    \hline
    Inscription avec email duplique & Erreur 400 & Valide \\
    \hline
    Connexion avec mot de passe incorrect & Erreur 401 & Valide \\
    \hline
    Connexion Google avec token invalide & Erreur 401 & Valide \\
    \hline
    Reinitialisation avec code expire & Erreur 400 & Valide \\
    \hline
    Consultation profil avec token invalide & Erreur 401 & Valide \\
    \hline
    Modification role par non-admin & Erreur 403 & Valide \\
    \hline
  \end{tabular}
\end{table}

\begin{table}[H]
  \centering
  \caption{Bilan du Sprint 1}
  \label{tab:bilan_s1}
  \begin{tabular}{|p{5cm}|c|p{6.5cm}|}
    \hline
    \textbf{Fonctionnalite} & \textbf{Statut} & \textbf{Observations} \\
    \hline
    Inscription locale & Fait & Validation MDP + unicite \\
    \hline
    Connexion locale & Fait & Controle erreurs et rate limiting \\
    \hline
    Connexion Google OAuth & Fait & Verification token + liaison compte \\
    \hline
    Mot de passe oublie/reinitialisation & Fait & Code temporaire + expiration \\
    \hline
    Gestion du profil & Fait & Consultation/modification + changement MDP \\
    \hline
    Administration utilisateurs & Fait & CRUD + roles + statuts \\
    \hline
  \end{tabular}
\end{table}

\section*{Conclusion}
\addcontentsline{toc}{section}{Conclusion}

Le Sprint 1 a permis de livrer un socle applicatif securise, stable et evolutif :
authentification locale et sociale, gestion des comptes, controle d'acces par roles,
et operations d'administration des utilisateurs.

\clearpage

% ##########################################################
%  CHAPITRE 3 – Sprint 2 : Contenus pedagogiques et QCM classique
% ##########################################################
\chapter{Sprint 2 : Contenus pedagogiques et QCM classique}
\minitoc
\newpage

\section*{Introduction}
\addcontentsline{toc}{section}{Introduction}

Le Sprint 2 vise la mise en oeuvre des contenus pedagogiques et des QCM classiques.
Il couvre deux axes principaux :

\begin{enumerate}
  \item La gestion complete des cours (creation, mise a jour, publication, consultation).
  \item Le module QCM classique (categories, creation, publication, passage et resultats).
\end{enumerate}

\section{Objectif et perimetre}

\subsection{Objectif general}

Fournir un module de cours fiable et un module QCM classique administrable, afin
de proposer un parcours d'apprentissage structure (cours + evaluation) avant
l'introduction des briques intelligentes du Sprint 3.

\subsection{Perimetre fonctionnel retenu}

\begin{enumerate}
  \item \textbf{Cours (Formateur/Admin)}
  \begin{itemize}
    \item Creer, modifier, supprimer un cours
    \item Publier / depublier un cours
    \item Ordonner et categoriser les cours
  \end{itemize}
  \item \textbf{Cours (Apprenant)}
  \begin{itemize}
    \item Consulter la liste des cours publies
    \item Consulter le detail d'un cours publie
  \end{itemize}
  \item \textbf{QCM classique (Admin)}
  \begin{itemize}
    \item Gerer les categories de QCM
    \item Creer, modifier, supprimer un QCM
    \item Publier / depublier un QCM
  \end{itemize}
  \item \textbf{QCM classique (Apprenant)}
  \begin{itemize}
    \item Consulter les QCM publies
    \item Passer un QCM et obtenir un score avec details
    \item Consulter l'historique des resultats
  \end{itemize}
\end{enumerate}

\section{Backlog Sprint 2}

\begin{table}[H]
  \centering
  \caption{Backlog du Sprint 2}
  \label{tab:backlog_s2}
  \begin{tabular}{|c|p{7.5cm}|c|c|}
    \hline
    \textbf{ID} & \textbf{User Story} & \textbf{Priorite} & \textbf{Est.} \\
    \hline
    S2-1 & En tant que formateur, je veux creer et modifier des cours & Haute & 8 SP \\
    \hline
    S2-2 & En tant que formateur, je veux publier ou depublier un cours & Haute & 3 SP \\
    \hline
    S2-3 & En tant qu'apprenant, je veux consulter les cours publies & Haute & 5 SP \\
    \hline
    S2-4 & En tant qu'admin, je veux gerer les categories de QCM & Moyenne & 5 SP \\
    \hline
    S2-5 & En tant qu'admin, je veux creer et gerer des QCM classiques & Haute & 8 SP \\
    \hline
    S2-6 & En tant qu'apprenant, je veux passer un QCM et voir mon score & Haute & 8 SP \\
    \hline
  \end{tabular}
\end{table}

\section{Analyse}

\subsection{Acteurs du Sprint 2}
\begin{enumerate}
  \item Formateur
  \item Administrateur
  \item Apprenant
\end{enumerate}

\subsection{Diagramme de cas d'utilisation -- Sprint 2}

\begin{figure}[H]
  \centering
  \resizebox{0.96\textwidth}{!}{%
  \begin{tikzpicture}[
    >=Stealth,
    uc/.style={draw, ellipse, minimum width=4.8cm, minimum height=0.95cm, align=center, font=\footnotesize},
    actor/.style={font=\small\bfseries, align=center},
    rel/.style={draw, ->, line width=0.6pt}
  ]
    \draw[dashed, rounded corners=6pt] (0,0) rectangle (18,-14.2);
    \node[font=\small\bfseries] at (9,-0.5)
      {Systeme -- Sprint 2 : Contenus pedagogiques et QCM classique};

    \node[uc] (UC1)  at (9,-2.0)  {Creer / modifier cours};
    \node[uc] (UC2)  at (9,-3.6)  {Publier / depublier cours};
    \node[uc] (UC3)  at (9,-5.2)  {Consulter cours publies};
    \node[uc] (UC4)  at (9,-6.8)  {Gerer categories QCM};
    \node[uc] (UC5)  at (9,-8.4)  {Creer / modifier QCM classique};
    \node[uc] (UC6)  at (9,-10.0) {Publier / depublier QCM};
    \node[uc] (UC7)  at (9,-11.6) {Passer un QCM et obtenir un score};
    \node[uc] (UC8)  at (9,-13.2) {Consulter historique des resultats};

    \node[actor] (For) at (-1.6,-3.2) {Formateur};
    \node[actor] (Adm) at (-1.6,-8.6) {Administrateur};
    \node[actor] (App) at (-1.6,-12.4) {Apprenant};

    \draw[rel] (For) -- (UC1);
    \draw[rel] (For) -- (UC2);
    \draw[rel] (For) -- (UC3);

    \draw[rel] (Adm) -- (UC1);
    \draw[rel] (Adm) -- (UC2);
    \draw[rel] (Adm) -- (UC4);
    \draw[rel] (Adm) -- (UC5);
    \draw[rel] (Adm) -- (UC6);

    \draw[rel] (App) -- (UC3);
    \draw[rel] (App) -- (UC7);
    \draw[rel] (App) -- (UC8);
  \end{tikzpicture}%
  }
  \caption{Diagramme de cas d'utilisation -- Sprint 2}
  \label{fig:uc_sprint2}
\end{figure}

\section{Conception}

\subsection{Cas d'utilisation raffines}

% ---- UC1 ----
\subsubsection{UC1 -- Creer ou modifier un cours}
\begin{table}[H]
  \centering
  \caption{Description du cas d'utilisation UC1}
  \begin{tabularx}{\textwidth}{>{\bfseries}p{4cm} X}
    \toprule
    Cas d'utilisation & Creer / modifier un cours \\
    \midrule
    Acteur principal & Formateur / Administrateur \\
    Preconditions & Compte authentifie avec role formateur ou admin. \\
    Scenario nominal & Saisie des metadonnees (titre, categorie, niveau), contenu du cours,
                       validation et enregistrement en base. \\
    Postconditions & Cours cree ou mis a jour et disponible en gestion. \\
    Exceptions & Champs invalides, cours introuvable (modification). \\
    \bottomrule
  \end{tabularx}
\end{table}

\ucseq{UC1 -- Creer / modifier cours}{Formateur}{API Courses}{CourseService + Base}{Saisie formulaire cours}{POST/PUT /api/courses}{Champs invalides ou cours introuvable}{201/200 Cours sauvegarde}

% ---- UC2 ----
\subsubsection{UC2 -- Publier / depublier un cours}
\begin{table}[H]
  \centering
  \caption{Description du cas d'utilisation UC2}
  \begin{tabularx}{\textwidth}{>{\bfseries}p{4cm} X}
    \toprule
    Cas d'utilisation & Publier / depublier un cours \\
    \midrule
    Acteur principal & Formateur / Administrateur \\
    Preconditions & Cours existant et droits de gestion. \\
    Scenario nominal & Changement du statut \texttt{is\_published} via endpoint de publication. \\
    Postconditions & Le cours est visible ou masque dans la liste publique. \\
    Exceptions & Cours introuvable, acces non autorise. \\
    \bottomrule
  \end{tabularx}
\end{table}

% ---- UC3 ----
\subsubsection{UC3 -- Consulter un cours publie}
\begin{table}[H]
  \centering
  \caption{Description du cas d'utilisation UC3}
  \begin{tabularx}{\textwidth}{>{\bfseries}p{4cm} X}
    \toprule
    Cas d'utilisation & Consulter un cours publie \\
    \midrule
    Acteur principal & Apprenant \\
    Preconditions & Cours publie disponible. \\
    Scenario nominal & Consultation de la liste, puis affichage detaille du cours. \\
    Postconditions & Contenu affiche avec metadonnees et ressources. \\
    Exceptions & Cours introuvable ou non publie. \\
    \bottomrule
  \end{tabularx}
\end{table}

% ---- UC4 ----
\subsubsection{UC4 -- Gerer les categories QCM}
\begin{table}[H]
  \centering
  \caption{Description du cas d'utilisation UC4}
  \begin{tabularx}{\textwidth}{>{\bfseries}p{4cm} X}
    \toprule
    Cas d'utilisation & Gerer les categories QCM \\
    \midrule
    Acteur principal & Administrateur \\
    Preconditions & Role admin authentifie. \\
    Scenario nominal & Creation, modification, activation/desactivation des categories. \\
    Postconditions & Categories disponibles pour la creation des QCM. \\
    Exceptions & Nom invalide, categorie utilisee ou deja existante. \\
    \bottomrule
  \end{tabularx}
\end{table}

% ---- UC5 ----
\subsubsection{UC5 -- Creer un QCM classique}
\begin{table}[H]
  \centering
  \caption{Description du cas d'utilisation UC5}
  \begin{tabularx}{\textwidth}{>{\bfseries}p{4cm} X}
    \toprule
    Cas d'utilisation & Creer un QCM classique \\
    \midrule
    Acteur principal & Administrateur \\
    Preconditions & Role admin, categorie active selectionnee. \\
    Scenario nominal & Saisie du QCM (titre, difficulte, duree, questions/reponses),
                       validation des reponses correctes, sauvegarde. \\
    Postconditions & QCM enregistre, pret a etre publie. \\
    Exceptions & Categorie invalide, question sans bonne reponse. \\
    \bottomrule
  \end{tabularx}
\end{table}

% ---- UC6 ----
\subsubsection{UC6 -- Passer un QCM et obtenir un score}
\begin{table}[H]
  \centering
  \caption{Description du cas d'utilisation UC6}
  \begin{tabularx}{\textwidth}{>{\bfseries}p{4cm} X}
    \toprule
    Cas d'utilisation & Passer un QCM et obtenir un score \\
    \midrule
    Acteur principal & Apprenant \\
    Preconditions & QCM publie disponible. \\
    Scenario nominal & Recuperation du QCM, soumission des reponses,
                       calcul du score et enregistrement du resultat. \\
    Postconditions & Score et details disponibles dans l'historique. \\
    Exceptions & QCM introuvable ou non publie. \\
    \bottomrule
  \end{tabularx}
\end{table}

\section{Diagrammes de sequence detailles du Sprint 2}

\subsection{Sequence detaillee -- Creation d'un cours}

\begin{figure}[H]
  \centering
  \begin{tikzpicture}[x=1pt, y=1pt,
    box/.style={draw, fill=gray!10, minimum width=1.8cm, minimum height=0.6cm, font=\footnotesize},
    msg/.style={draw, ->, >=Stealth, font=\scriptsize},
    ret/.style={draw, dashed, ->, >=Stealth, font=\scriptsize},
    life/.style={draw, dashed, thin}
  ]
  \foreach \x/\lbl in
    {0/Formateur, 90/Frontend, 190/API Courses, 300/Service, 410/Base MySQL}{
    \node[box] at (\x,0) {\lbl};
    \draw[life] (\x,-8) -- (\x,-300);
  }
  \draw[msg] (0,-30)  -- node[above]{donnees cours} (90,-30);
  \draw[msg] (90,-55) -- node[above]{POST /api/courses} (190,-55);
  \draw[msg] (190,-80) -- node[above]{controle role} (300,-80);
  \draw[ret] (300,-100) -- node[above]{OK} (190,-100);
  \draw[msg] (190,-125) -- node[above]{insert cours} (410,-125);
  \draw[ret] (410,-145) -- node[above]{course} (190,-145);
  \draw[ret] (190,-170) -- node[above]{201 CourseDetail} (90,-170);
  \draw[ret] (90,-190)  -- node[above]{cours cree} (0,-190);
  \end{tikzpicture}
  \caption{Diagramme de sequence detaille -- Creation d'un cours}
  \label{fig:seq_create_course}
\end{figure}

\subsection{Sequence detaillee -- Passage d'un QCM}

\begin{figure}[H]
  \centering
  \begin{tikzpicture}[x=1pt, y=1pt,
    box/.style={draw, fill=gray!10, minimum width=1.8cm, minimum height=0.6cm, font=\footnotesize},
    msg/.style={draw, ->, >=Stealth, font=\scriptsize},
    ret/.style={draw, dashed, ->, >=Stealth, font=\scriptsize},
    life/.style={draw, dashed, thin}
  ]
  \foreach \x/\lbl in
    {0/Apprenant, 90/Frontend, 190/API QCM, 300/Service, 410/Base MySQL}{
    \node[box] at (\x,0) {\lbl};
    \draw[life] (\x,-8) -- (\x,-320);
  }
  \draw[msg] (0,-30)  -- node[above]{selection QCM} (90,-30);
  \draw[msg] (90,-55) -- node[above]{GET /api/qcm/\{id\}} (190,-55);
  \draw[msg] (190,-80) -- node[above]{charger QCM publie} (410,-80);
  \draw[ret] (410,-100) -- node[above]{QCM + questions} (190,-100);
  \draw[ret] (190,-120) -- node[above]{200 QCM public} (90,-120);

  \draw[msg] (0,-150)  -- node[above]{soumission reponses} (90,-150);
  \draw[msg] (90,-175) -- node[above]{POST /api/qcm/\{id\}/submit} (190,-175);
  \draw[msg] (190,-200) -- node[above]{calcul score + sauvegarde} (300,-200);
  \draw[msg] (300,-225) -- node[above]{insert resultat} (410,-225);
  \draw[ret] (410,-245) -- node[above]{OK} (190,-245);
  \draw[ret] (190,-270) -- node[above]{200 QCMResult} (90,-270);
  \draw[ret] (90,-290)  -- node[above]{affichage score} (0,-290);
  \end{tikzpicture}
  \caption{Diagramme de sequence detaille -- Passage d'un QCM}
  \label{fig:seq_submit_qcm}
\end{figure}

\section{Diagramme de classes -- Sprint 2}

\begin{figure}[H]
  \centering
  \begin{tikzpicture}[
    cls/.style={draw, rounded corners=2pt, fill=gray!10, align=left,
                font=\scriptsize, text width=5.2cm, inner sep=4pt},
    assoc/.style={draw, ->, >=Stealth, line width=0.6pt}
  ]
  \node[cls] (User) at (-4.2,0) {
    \textbf{User}\\
    \rule{4.8cm}{0.3pt}\\
    id: int\\
    username: string\\
    email: string\\
    role: UserRole\\
    is\_active: bool
  };

  \node[cls] (Course) at (0,0) {
    \textbf{Course}\\
    \rule{4.8cm}{0.3pt}\\
    id: int\\
    title: string\\
    category: string\\
    level: enum\\
    content: text\\
    is\_published: bool\\
    created\_by: int
  };

  \node[cls] (QCMCat) at (8.6,0) {
    \textbf{QCMCategory}\\
    \rule{4.8cm}{0.3pt}\\
    id: int\\
    slug: string\\
    name: string\\
    is\_active: bool
  };

  \node[cls] (QCM) at (0,-5.8) {
    \textbf{QCM}\\
    \rule{4.8cm}{0.3pt}\\
    id: int\\
    title: string\\
    description: text\\
    category: string\\
    difficulty: string\\
    duration\_minutes: int\\
    pass\_score: int\\
    is\_published: bool\\
    created\_by: int
  };

  \node[cls] (Question) at (8.6,-5.8) {
    \textbf{Question}\\
    \rule{4.8cm}{0.3pt}\\
    id: int\\
    qcm\_id: int\\
    text: text\\
    image\_url: string\\
    explanation: text\\
    order: int
  };

  \node[cls] (Answer) at (8.6,-11.2) {
    \textbf{Answer}\\
    \rule{4.8cm}{0.3pt}\\
    id: int\\
    question\_id: int\\
    text: text\\
    is\_correct: bool
  };

  \node[cls] (Result) at (0,-11.2) {
    \textbf{UserQCMResult}\\
    \rule{4.8cm}{0.3pt}\\
    id: int\\
    user\_id: int\\
    qcm\_id: int\\
    score: float\\
    total\_questions: int\\
    correct\_answers: int\\
    passed: bool\\
    completed\_at: datetime
  };

  \draw[assoc] (User.east) -- node[above]{cree} (Course.west);
  \draw[assoc] (User.south) -- node[left]{cree} (QCM.north west);
  \draw[assoc] (User.south east) -- node[left]{obtient} (Result.north west);
  \draw[assoc] (QCMCat.south) -- node[right]{categorie} (QCM.north);
  \draw[assoc] (QCM.east) -- node[above]{contient} (Question.west);
  \draw[assoc] (Question.south) -- node[right]{contient} (Answer.north);
  \draw[assoc] (QCM.south) -- node[left]{resultats} (Result.north);
  \end{tikzpicture}
  \caption{Diagramme de classes -- Sprint 2}
  \label{fig:class_sprint2}
\end{figure}

\section{Extraits de code -- Sprint 2}

\subsection{API Cours (courses.py)}

\begin{lstlisting}[language=Python,
  caption={courses.py -- Creation, consultation et publication des cours}]
@router.post("/", response_model=CourseDetailResponse, status_code=201)
def create_course(data: CourseCreate,
                  user: User = Depends(require_formateur_or_admin),
                  db: Session = Depends(get_db)):
    course = Course(
        title=data.title, description=data.description,
        category=data.category, level=data.level,
        duration=data.duration, image_url=data.image_url,
        content=data.content, video_url=data.video_url,
        order=data.order, created_by=user.id,
    )
    db.add(course)
    db.commit()
    db.refresh(course)
    return CourseDetailResponse.model_validate(course)


@router.get("/published", response_model=List[CourseListResponse])
def list_published_courses(category: Optional[str] = None,
                           level: Optional[str] = None,
                           db: Session = Depends(get_db)):
    query = db.query(Course).filter(Course.is_published == True)
    if category:
        query = query.filter(Course.category == category)
    if level:
        query = query.filter(Course.level == level)
    return [CourseListResponse.model_validate(c) for c in query.all()]


@router.put("/{course_id}/publish")
def toggle_publish(course_id: int,
                   user: User = Depends(require_formateur_or_admin),
                   db: Session = Depends(get_db)):
    course = db.query(Course).filter(Course.id == course_id).first()
    if not course:
        raise HTTPException(404, "Cours non trouve")
    course.is_published = not course.is_published
    db.commit()
    return {"is_published": course.is_published}
\end{lstlisting}

\subsection{API QCM (qcm.py)}

\begin{lstlisting}[language=Python,
  caption={qcm.py -- Categories, creation et soumission de QCM classiques}]
@router.post("/categories", response_model=QCMCategoryResponse, status_code=201)
def create_category(data: QCMCategoryCreate,
                    admin: User = Depends(require_admin),
                    db: Session = Depends(get_db)):
    slug = _slugify(data.name)
    row = QCMCategory(slug=slug, name=data.name.strip(), is_active=True)
    db.add(row)
    db.commit()
    db.refresh(row)
    return row


@router.post("/", response_model=QCMDetailResponse, status_code=201)
def create_qcm(data: QCMCreate,
               admin: User = Depends(require_admin),
               db: Session = Depends(get_db)):
    category = _require_valid_active_category(db, data.category)
    qcm = QCM(title=data.title, description=data.description,
              category=category.slug, difficulty=data.difficulty,
              duration_minutes=data.duration_minutes,
              pass_score=data.pass_score, created_by=admin.id)
    db.add(qcm)
    db.flush()
    for q_data in data.questions:
        question = Question(qcm_id=qcm.id, text=q_data.text,
                            image_url=q_data.image_url,
                            explanation=q_data.explanation,
                            order=q_data.order)
        db.add(question)
        db.flush()
        for a_data in q_data.answers:
            db.add(Answer(question_id=question.id,
                          text=a_data.text, is_correct=a_data.is_correct))
    db.commit()
    return _qcm_to_detail_response(qcm, _category_name_cache(db))


@router.post("/{qcm_id}/submit", response_model=QCMResultResponse)
def submit_qcm(qcm_id: int, data: SubmitQCMRequest,
               current_user: User = Depends(get_current_user),
               db: Session = Depends(get_db)):
    qcm = db.query(QCM).filter(QCM.id == qcm_id,
                               QCM.is_published == True).first()
    if not qcm:
        raise HTTPException(404, "QCM non trouve")
    # calcul score + sauvegarde resultat
\end{lstlisting}

\section{Captures d'ecran des interfaces principales}

\begin{figure}[H]
  \centering
  \IfFileExists{captures/courses_list.png}{%
    \includegraphics[width=0.90\textwidth]{captures/courses_list.png}%
  }{%
    \begin{tikzpicture}
      \fill[gray!12] (0,0) rectangle (13.5cm,-6.2cm);
      \draw[gray!40] (0,0) rectangle (13.5cm,-6.2cm);
      \node[font=\small\itshape] at (6.75cm,-3.1cm) {[Capture liste des cours]};
    \end{tikzpicture}%
  }
  \caption{Interface de consultation des cours}
\end{figure}

\begin{figure}[H]
  \centering
  \IfFileExists{captures/qcm_list.png}{%
    \includegraphics[width=0.90\textwidth]{captures/qcm_list.png}%
  }{%
    \begin{tikzpicture}
      \fill[gray!12] (0,0) rectangle (13.5cm,-6.2cm);
      \draw[gray!40] (0,0) rectangle (13.5cm,-6.2cm);
      \node[font=\small\itshape] at (6.75cm,-3.1cm) {[Capture liste QCM]};
    \end{tikzpicture}%
  }
  \caption{Interface des QCM classiques}
\end{figure}

\section{Evaluation du Sprint 2}

\subsection{Scenarios de tests realises}

\begin{table}[H]
  \centering
  \caption{Table de tests fonctionnels du Sprint 2}
  \label{tab:tests_s2}
  \begin{tabular}{|p{5.2cm}|p{3.2cm}|p{5.4cm}|}
    \hline
    \textbf{Scenario} & \textbf{Resultat attendu} & \textbf{Statut} \\
    \hline
    Creation cours avec donnees valides & 201 + cours cree & Valide \\
    \hline
    Publication cours inexistant & Erreur 404 & Valide \\
    \hline
    Consultation cours publies & Liste non vide & Valide \\
    \hline
    Creation categorie QCM & 201 + categorie creee & Valide \\
    \hline
    Creation QCM sans bonne reponse & Erreur 400 & Valide \\
    \hline
    Passage QCM publie & Score calcule + resultat sauvegarde & Valide \\
    \hline
  \end{tabular}
\end{table}

\begin{table}[H]
  \centering
  \caption{Bilan du Sprint 2}
  \label{tab:bilan_s2}
  \begin{tabular}{|p{5cm}|c|p{6.5cm}|}
    \hline
    \textbf{Fonctionnalite} & \textbf{Statut} & \textbf{Observations} \\
    \hline
    CRUD cours & Fait & Creation, mise a jour, suppression \\
    \hline
    Publication cours & Fait & Controle par role formateur/admin \\
    \hline
    Consultation cours publies & Fait & Filtre categorie/niveau \\
    \hline
    Gestion categories QCM & Fait & Activation/desactivation + contraintes \\
    \hline
    QCM classique (admin) & Fait & Questions/reponses valides \\
    \hline
    Passage QCM (apprenant) & Fait & Score, details, historique \\
    \hline
  \end{tabular}
\end{table}

\section*{Conclusion}
\addcontentsline{toc}{section}{Conclusion}

Le Sprint 2 a mis en place les contenus pedagogiques et le module QCM classique.
Les cours sont geres par les formateurs et l'administration, tandis que les QCM
peuvent etre crees, publies et passes par les apprenants avec calcul du score
et sauvegarde de l'historique.

\clearpage

% ##########################################################
%  CHAPITRE 4 – Sprint 3 : Systemes intelligents et evaluation
% ##########################################################
\chapter{Sprint 3 : Systemes intelligents et evaluation}
\minitoc
\newpage

\section*{Introduction}
\addcontentsline{toc}{section}{Introduction}

Le Sprint 3 introduit les composants intelligents du systeme : generation
automatique de QCM par LLM, assistant conversationnel RAG, ingestion des
documents pedagogiques et detection de panneaux par image.

\section{Objectif et perimetre}

\subsection{Objectif general}

Fournir une experience d'apprentissage personnalisee a travers des services IA
fiables, avec un moteur RAG cite et une generation de QCM adaptee au niveau
et au theme.

\subsection{Perimetre fonctionnel retenu}

\begin{enumerate}
  \item \textbf{QCM intelligent (Apprenant)}
  \begin{itemize}
    \item Generer un QCM par LLM (general ou par theme)
    \item Configurer le nombre de questions, duree et difficulte
    \item Passer le QCM genere et obtenir un score detaille
  \end{itemize}
  \item \textbf{Chatbot RAG (Apprenant)}
  \begin{itemize}
    \item Poser une question en FR/AR
    \item Obtenir une reponse citee avec sources
    \item Donner un feedback sur la reponse
  \end{itemize}
  \item \textbf{Administration IA (Admin)}
  \begin{itemize}
    \item Lancer l'ingestion documentaire et suivre le vector store
    \item Vider le store si necessaire
  \end{itemize}
  \item \textbf{Vision par ordinateur}
  \begin{itemize}
    \item Detecter les panneaux de signalisation a partir d'une image
  \end{itemize}
\end{enumerate}

\section{Backlog Sprint 3}

\begin{table}[H]
  \centering
  \caption{Backlog du Sprint 3}
  \label{tab:backlog_s3}
  \begin{tabular}{|c|p{7.5cm}|c|c|}
    \hline
    \textbf{ID} & \textbf{User Story} & \textbf{Priorite} & \textbf{Est.} \\
    \hline
    S3-1 & En tant qu'apprenant, je veux generer un QCM automatique par theme & Haute & 8 SP \\
    \hline
    S3-2 & En tant qu'apprenant, je veux configurer mon QCM (duree, difficulte) & Haute & 5 SP \\
    \hline
    S3-3 & En tant qu'apprenant, je veux discuter avec un chatbot RAG & Haute & 8 SP \\
    \hline
    S3-4 & En tant qu'utilisateur, je veux donner un feedback sur la reponse du chatbot & Moyenne & 3 SP \\
    \hline
    S3-5 & En tant qu'admin, je veux ingerer les documents pedagogiques & Haute & 5 SP \\
    \hline
    S3-6 & En tant qu'utilisateur, je veux detecter un panneau par image & Moyenne & 8 SP \\
    \hline
  \end{tabular}
\end{table}

\section{Analyse}

\subsection{Acteurs du Sprint 3}
\begin{enumerate}
  \item Apprenant
  \item Administrateur
  \item Services IA (LLM, RAG, Vision)
  \item Vector store (embeddings)
\end{enumerate}

\subsection{Diagramme de cas d'utilisation -- Sprint 3}

\begin{figure}[H]
  \centering
  \resizebox{0.96\textwidth}{!}{%
  \begin{tikzpicture}[
    >=Stealth,
    uc/.style={draw, ellipse, minimum width=4.9cm, minimum height=0.95cm, align=center, font=\footnotesize},
    actor/.style={font=\small\bfseries, align=center},
    rel/.style={draw, ->, line width=0.6pt}
  ]
    \draw[dashed, rounded corners=6pt] (0,0) rectangle (18,-14.6);
    \node[font=\small\bfseries] at (9,-0.5)
      {Systeme -- Sprint 3 : IA, RAG et evaluation};

    \node[uc] (UC1) at (9,-2.0)  {Generer QCM intelligent (LLM)};
    \node[uc] (UC2) at (9,-3.6)  {Configurer QCM (duree, difficulte)};
    \node[uc] (UC3) at (9,-5.2)  {Passer QCM genere};
    \node[uc] (UC4) at (9,-6.8)  {Discuter avec le chatbot (RAG)};
    \node[uc] (UC5) at (9,-8.4)  {Donner un feedback};
    \node[uc] (UC6) at (9,-10.0) {Ingerer documents pedagogiques};
    \node[uc] (UC7) at (9,-11.6) {Consulter stats vector store};
    \node[uc] (UC8) at (9,-13.2) {Detecter panneau par image};

    \node[actor] (App) at (-1.6,-6.0) {Apprenant};
    \node[actor] (Adm) at (-1.6,-11.0) {Administrateur};

    \draw[rel] (App) -- (UC1);
    \draw[rel] (App) -- (UC2);
    \draw[rel] (App) -- (UC3);
    \draw[rel] (App) -- (UC4);
    \draw[rel] (App) -- (UC5);
    \draw[rel] (App) -- (UC8);

    \draw[rel] (Adm) -- (UC6);
    \draw[rel] (Adm) -- (UC7);
  \end{tikzpicture}%
  }
  \caption{Diagramme de cas d'utilisation -- Sprint 3}
  \label{fig:uc_sprint3}
\end{figure}

\section{Conception}

\subsection{Cas d'utilisation raffines}

\subsubsection{UC1 -- Generer un QCM intelligent}
\begin{table}[H]
  \centering
  \caption{Description du cas d'utilisation UC1}
  \begin{tabularx}{\textwidth}{>{\bfseries}p{4cm} X}
    \toprule
    Cas d'utilisation & Generer un QCM intelligent \\
    \midrule
    Acteur principal & Apprenant \\
    Preconditions & Utilisateur authentifie, parametres valides. \\
    Scenario nominal & Envoi du theme et des parametres, generation LLM,
                       sauvegarde du QCM genere. \\
    Postconditions & QCM genere pret a etre passe. \\
    Exceptions & Theme manquant en mode specifique, erreur IA. \\
    \bottomrule
  \end{tabularx}
\end{table}

\subsubsection{UC2 -- Discuter avec le chatbot (RAG)}
\begin{table}[H]
  \centering
  \caption{Description du cas d'utilisation UC2}
  \begin{tabularx}{\textwidth}{>{\bfseries}p{4cm} X}
    \toprule
    Cas d'utilisation & Discuter avec le chatbot (RAG) \\
    \midrule
    Acteur principal & Apprenant \\
    Preconditions & Documents ingeress et indexations disponibles. \\
    Scenario nominal & Recherche semantique, generation reponse citee,
                       sauvegarde conversation si authentifie. \\
    Postconditions & Reponse citee retournee. \\
    Exceptions & Aucun document pertinent, erreur de generation. \\
    \bottomrule
  \end{tabularx}
\end{table}

\subsubsection{UC3 -- Ingerer les documents}
\begin{table}[H]
  \centering
  \caption{Description du cas d'utilisation UC3}
  \begin{tabularx}{\textwidth}{>{\bfseries}p{4cm} X}
    \toprule
    Cas d'utilisation & Ingerer les documents pedagogiques \\
    \midrule
    Acteur principal & Administrateur \\
    Preconditions & Documents disponibles dans les dossiers fr/ar. \\
    Scenario nominal & Chargement, decoupage, vectorisation, indexation. \\
    Postconditions & Vector store alimente et disponible pour le RAG. \\
    Exceptions & Aucun document ou erreur d'indexation. \\
    \bottomrule
  \end{tabularx}
\end{table}

\subsubsection{UC4 -- Detecter un panneau par image}
\begin{table}[H]
  \centering
  \caption{Description du cas d'utilisation UC4}
  \begin{tabularx}{\textwidth}{>{\bfseries}p{4cm} X}
    \toprule
    Cas d'utilisation & Detecter un panneau par image \\
    \midrule
    Acteur principal & Apprenant \\
    Preconditions & Image valide (format, taille). \\
    Scenario nominal & Analyse vision, matching base panneaux, retour du resultat. \\
    Postconditions & Panneaux detectes affiches a l'utilisateur. \\
    Exceptions & Format non supporte, aucune detection. \\
    \bottomrule
  \end{tabularx}
\end{table}

\section{Diagrammes de sequence detailles du Sprint 3}

\subsection{Sequence detaillee -- Generation d'un QCM par LLM}

\begin{figure}[H]
  \centering
  \begin{tikzpicture}[x=1pt, y=1pt,
    box/.style={draw, fill=gray!10, minimum width=1.8cm, minimum height=0.6cm, font=\footnotesize},
    msg/.style={draw, ->, >=Stealth, font=\scriptsize},
    ret/.style={draw, dashed, ->, >=Stealth, font=\scriptsize},
    life/.style={draw, dashed, thin}
  ]
  \foreach \x/\lbl in
    {0/Apprenant, 90/Frontend, 190/API QCM, 300/QCMGenerator, 410/LLM}{
    \node[box] at (\x,0) {\lbl};
    \draw[life] (\x,-8) -- (\x,-300);
  }
  \draw[msg] (0,-30)  -- node[above]{parametres QCM} (90,-30);
  \draw[msg] (90,-55) -- node[above]{POST /api/qcm/generate} (190,-55);
  \draw[msg] (190,-80) -- node[above]{prepare prompt} (300,-80);
  \draw[msg] (300,-105) -- node[above]{generation} (410,-105);
  \draw[ret] (410,-125) -- node[above]{questions JSON} (300,-125);
  \draw[msg] (300,-150) -- node[above]{sauvegarde QCM} (190,-150);
  \draw[ret] (190,-175) -- node[above]{201 QCMGenerate} (90,-175);
  \draw[ret] (90,-195)  -- node[above]{QCM genere} (0,-195);
  \end{tikzpicture}
  \caption{Diagramme de sequence detaille -- Generation QCM intelligent}
  \label{fig:seq_gen_qcm}
\end{figure}

\subsection{Sequence detaillee -- Question au chatbot RAG}

\begin{figure}[H]
  \centering
  \begin{tikzpicture}[x=1pt, y=1pt,
    box/.style={draw, fill=gray!10, minimum width=1.8cm, minimum height=0.6cm, font=\footnotesize},
    msg/.style={draw, ->, >=Stealth, font=\scriptsize},
    ret/.style={draw, dashed, ->, >=Stealth, font=\scriptsize},
    life/.style={draw, dashed, thin}
  ]
  \foreach \x/\lbl in
    {0/Apprenant, 90/Frontend, 190/API Chat, 300/RAGService, 410/VectorStore}{
    \node[box] at (\x,0) {\lbl};
    \draw[life] (\x,-8) -- (\x,-330);
  }
  \draw[msg] (0,-30)  -- node[above]{question} (90,-30);
  \draw[msg] (90,-55) -- node[above]{POST /api/chat/ask} (190,-55);
  \draw[msg] (190,-80) -- node[above]{search\_similar()} (410,-80);
  \draw[ret] (410,-100) -- node[above]{chunks pertinents} (300,-100);
  \draw[msg] (300,-125) -- node[above]{LLM + citations} (190,-125);
  \draw[ret] (190,-150) -- node[above]{200 ChatResponse} (90,-150);
  \draw[ret] (90,-170)  -- node[above]{reponse affichee} (0,-170);
  \end{tikzpicture}
  \caption{Diagramme de sequence detaille -- Question au chatbot RAG}
  \label{fig:seq_chat_rag}
\end{figure}

\section{Diagramme de classes -- Sprint 3}

\begin{figure}[H]
  \centering
  \begin{tikzpicture}[
    cls/.style={draw, rounded corners=2pt, fill=gray!10, align=left,
                font=\scriptsize, text width=5.2cm, inner sep=4pt},
    assoc/.style={draw, ->, >=Stealth, line width=0.6pt}
  ]
  \node[cls] (User) at (-4.2,0) {
    \textbf{User}\\
    \rule{4.8cm}{0.3pt}\\
    id: int\\
    username: string\\
    role: UserRole
  };

  \node[cls] (Conversation) at (0,0) {
    \textbf{Conversation}\\
    \rule{4.8cm}{0.3pt}\\
    id: int\\
    user\_id: int\\
    title: string\\
    language: string
  };

  \node[cls] (Message) at (8.6,0) {
    \textbf{Message}\\
    \rule{4.8cm}{0.3pt}\\
    id: int\\
    conversation\_id: int\\
    role: string\\
    content: text\\
    language: string
  };

  \node[cls] (Feedback) at (8.6,-6.0) {
    \textbf{Feedback}\\
    \rule{4.8cm}{0.3pt}\\
    id: int\\
    session\_id: string\\
    question: text\\
    answer: text\\
    is\_positive: bool
  };

  \node[cls] (QCM) at (0,-6.0) {
    \textbf{QCM}\\
    \rule{4.8cm}{0.3pt}\\
    id: int\\
    title: string\\
    is\_generated: bool\\
    generation\_mode: string\\
    generation\_theme: string
  };

  \draw[assoc] (User.east) -- node[above]{possede} (Conversation.west);
  \draw[assoc] (Conversation.east) -- node[above]{contient} (Message.west);
  \draw[assoc] (User.south) -- node[left]{genere} (QCM.north);
  \draw[assoc] (User.south east) -- node[right]{donne} (Feedback.west);
  \end{tikzpicture}
  \caption{Diagramme de classes -- Sprint 3}
  \label{fig:class_sprint3}
\end{figure}

\section{Extraits de code -- Sprint 3}

\subsection{Generation QCM (qcm.py)}

\begin{lstlisting}[language=Python,
  caption={qcm.py -- Generation automatique de QCM par LLM}]
@router.post("/generate", response_model=QCMGenerateResponse)
def generate_qcm(data: QCMGenerateRequest,
                 current_user: User = Depends(get_current_user),
                 db: Session = Depends(get_db)):
    if data.question_count not in {20, 25, 30}:
        raise HTTPException(400, "Nombre de questions invalide")
    payload = generate_qcm_from_docs(
        language=data.language,
        mode=data.mode,
        theme=data.theme,
        question_count=data.question_count,
        difficulty=data.difficulty,
    )
    # creation QCM + questions
\end{lstlisting}

\subsection{Chatbot RAG (chat.py)}

\begin{lstlisting}[language=Python,
  caption={chat.py -- Endpoint de question au chatbot RAG}]
@router.post("/ask", response_model=ChatResponse)
async def ask(request: ChatRequest,
              conversation_id: Optional[int] = None,
              current_user: Optional[User] = Depends(get_current_user_optional),
              db: Session = Depends(get_db)):
    result = ask_question(
        question=request.question,
        language=request.language,
    )
    # sauvegarde optionnelle dans conversation
    return ChatResponse(**result)
\end{lstlisting}

\subsection{Ingestion RAG (chat.py)}

\begin{lstlisting}[language=Python,
  caption={chat.py -- Ingestion des documents dans le vector store}]
@router.post("/ingest", response_model=IngestResponse)
async def ingest_documents(request: IngestRequest):
    if request.clear_existing:
        clear_vector_store()
    documents = load_documents()
    chunks = split_documents(documents)
    add_documents_to_store(chunks)
    return IngestResponse(status="success")
\end{lstlisting}

\subsection{Detection de panneaux (image\_detector.py)}

\begin{lstlisting}[language=Python,
  caption={image_detector.py -- Pipeline de detection vision}]
def detect_sign_from_image(image_base64: str, language: str = "fr") -> dict:
    vision_response = analyze_image_with_vision(image_base64, language)
    matched_signs = match_vision_to_signs(vision_response, language)
    return {
        "signs": matched_signs,
        "description": build_description(matched_signs, vision_response, language),
    }
\end{lstlisting}

\section{Captures d'ecran des interfaces principales}

\begin{figure}[H]
  \centering
  \IfFileExists{captures/qcm_generate.png}{%
    \includegraphics[width=0.90\textwidth]{captures/qcm_generate.png}%
  }{%
    \begin{tikzpicture}
      \fill[gray!12] (0,0) rectangle (13.5cm,-6.2cm);
      \draw[gray!40] (0,0) rectangle (13.5cm,-6.2cm);
      \node[font=\small\itshape] at (6.75cm,-3.1cm) {[Capture generation QCM intelligent]};
    \end{tikzpicture}%
  }
  \caption{Interface de generation de QCM intelligent}
\end{figure}

\begin{figure}[H]
  \centering
  \IfFileExists{captures/chatbot.png}{%
    \includegraphics[width=0.90\textwidth]{captures/chatbot.png}%
  }{%
    \begin{tikzpicture}
      \fill[gray!12] (0,0) rectangle (13.5cm,-6.2cm);
      \draw[gray!40] (0,0) rectangle (13.5cm,-6.2cm);
      \node[font=\small\itshape] at (6.75cm,-3.1cm) {[Capture interface chatbot RAG]};
    \end{tikzpicture}%
  }
  \caption{Interface du chatbot intelligent}
\end{figure}

\begin{figure}[H]
  \centering
  \IfFileExists{captures/detect_sign.png}{%
    \includegraphics[width=0.90\textwidth]{captures/detect_sign.png}%
  }{%
    \begin{tikzpicture}
      \fill[gray!12] (0,0) rectangle (13.5cm,-6.2cm);
      \draw[gray!40] (0,0) rectangle (13.5cm,-6.2cm);
      \node[font=\small\itshape] at (6.75cm,-3.1cm) {[Capture detection de panneaux]};
    \end{tikzpicture}%
  }
  \caption{Interface de detection de panneaux}
\end{figure}

\section{Evaluation du Sprint 3}

\subsection{Scenarios de tests realises}

\begin{table}[H]
  \centering
  \caption{Table de tests fonctionnels du Sprint 3}
  \label{tab:tests_s3}
  \begin{tabular}{|p{5.2cm}|p{3.2cm}|p{5.4cm}|}
    \hline
    \textbf{Scenario} & \textbf{Resultat attendu} & \textbf{Statut} \\
    \hline
    Generation QCM general & QCM genere et sauvegarde & Valide \\
    \hline
    Generation QCM theme sans theme & Erreur 400 & Valide \\
    \hline
    Question chatbot sans docs & Message "aucune info" & Valide \\
    \hline
    Ingestion documents & Chunks indexes dans le store & Valide \\
    \hline
    Detection image non supportee & Erreur 400 & Valide \\
    \hline
  \end{tabular}
\end{table}

\begin{table}[H]
  \centering
  \caption{Bilan du Sprint 3}
  \label{tab:bilan_s3}
  \begin{tabular}{|p{5cm}|c|p{6.5cm}|}
    \hline
    \textbf{Fonctionnalite} & \textbf{Statut} & \textbf{Observations} \\
    \hline
    QCM intelligent & Fait & Generation auto + anti-repetition \\
    \hline
    Chatbot RAG & Fait & Reponses citees FR/AR \\
    \hline
    Ingestion documents & Fait & Decoupage + vectorisation \\
    \hline
    Feedback chatbot & Fait & Stats positives/negatives \\
    \hline
    Detection panneaux & Fait & Vision + matching base locale \\
    \hline
  \end{tabular}
\end{table}

\section*{Conclusion}
\addcontentsline{toc}{section}{Conclusion}

Le Sprint 3 a apporte l'intelligence applicative : generation de QCM par LLM,
assistant RAG bilingue avec citations, ingestion documentaire et detection de
panneaux par vision. Ces modules preparent l'etape finale de stabilisation et
d'optimisation.

\clearpage

% ##########################################################
%  CHAPITRE 5 – Sprint 4 : Finalisation et validation du systeme
% ##########################################################
\chapter{Sprint 4 : Finalisation et validation du systeme}
\minitoc
\newpage

\section*{Introduction}
\addcontentsline{toc}{section}{Introduction}

Le Sprint 4 est consacre a la stabilisation globale du produit, la correction
des anomalies, la validation fonctionnelle et la preparation a la livraison.
Il consolide l'integration frontend/backend, les tests de regression et la
documentation finale.

\section{Objectif et perimetre}

\subsection{Objectif general}

Valider l'ensemble des modules et garantir un produit deployable, robuste et
coherent avant la soutenance.

\subsection{Perimetre fonctionnel retenu}

\begin{enumerate}
  \item \textbf{Integration}
  \begin{itemize}
    \item Verification des flux complets (auth, cours, QCM, IA, stats)
    \item Harmonisation des interfaces et des messages d'erreur
  \end{itemize}
  \item \textbf{Validation et tests}
  \begin{itemize}
    \item Tests fonctionnels de regression
    \item Tests de securite de base (roles, droits, tokens)
    \item Tests de performance elementaires (temps de reponse)
  \end{itemize}
  \item \textbf{Deploiement et documentation}
  \begin{itemize}
    \item Preparation des scripts de demarrage
    \item Nettoyage des logs et donnees de test
    \item Finalisation de la documentation technique et utilisateur
  \end{itemize}
\end{enumerate}

\section{Backlog Sprint 4}

\begin{table}[H]
  \centering
  \caption{Backlog du Sprint 4}
  \label{tab:backlog_s4}
  \begin{tabular}{|c|p{7.5cm}|c|c|}
    \hline
    \textbf{ID} & \textbf{User Story} & \textbf{Priorite} & \textbf{Est.} \\
    \hline
    S4-1 & En tant qu'equipe, je veux valider tous les flux metier en bout en bout & Haute & 8 SP \\
    \hline
    S4-2 & En tant qu'equipe, je veux corriger les anomalies identifiees & Haute & 5 SP \\
    \hline
    S4-3 & En tant qu'equipe, je veux verifier la securite des acces & Haute & 5 SP \\
    \hline
    S4-4 & En tant qu'equipe, je veux preparer un deploiement stable & Moyenne & 5 SP \\
    \hline
    S4-5 & En tant qu'equipe, je veux finaliser la documentation technique & Moyenne & 3 SP \\
    \hline
    S4-6 & En tant qu'equipe, je veux preparer la demo de soutenance & Moyenne & 3 SP \\
    \hline
  \end{tabular}
\end{table}

\section{Analyse}

\subsection{Acteurs du Sprint 4}
\begin{enumerate}
  \item Equipe de developpement
  \item Administrateur
  \item Utilisateurs test (apprenants / formateurs)
\end{enumerate}

\subsection{Diagramme de cas d'utilisation -- Sprint 4}

\begin{figure}[H]
  \centering
  \resizebox{0.96\textwidth}{!}{%
  \begin{tikzpicture}[
    >=Stealth,
    uc/.style={draw, ellipse, minimum width=5.2cm, minimum height=0.95cm, align=center, font=\footnotesize},
    actor/.style={font=\small\bfseries, align=center},
    rel/.style={draw, ->, line width=0.6pt}
  ]
    \draw[dashed, rounded corners=6pt] (0,0) rectangle (18,-12.8);
    \node[font=\small\bfseries] at (9,-0.5)
      {Systeme -- Sprint 4 : Finalisation et validation};

    \node[uc] (UC1) at (9,-2.0)  {Executer tests de regression};
    \node[uc] (UC2) at (9,-3.8)  {Corriger anomalies};
    \node[uc] (UC3) at (9,-5.6)  {Verifier securite et roles};
    \node[uc] (UC4) at (9,-7.4)  {Valider performances};
    \node[uc] (UC5) at (9,-9.2)  {Deployer et stabiliser};
    \node[uc] (UC6) at (9,-11.0) {Finaliser documentation};

    \node[actor] (Dev) at (-1.6,-5.0) {Equipe dev};
    \node[actor] (Adm) at (-1.6,-10.0) {Administrateur};

    \draw[rel] (Dev) -- (UC1);
    \draw[rel] (Dev) -- (UC2);
    \draw[rel] (Dev) -- (UC3);
    \draw[rel] (Dev) -- (UC4);
    \draw[rel] (Dev) -- (UC5);
    \draw[rel] (Dev) -- (UC6);

    \draw[rel] (Adm) -- (UC5);
    \draw[rel] (Adm) -- (UC6);
  \end{tikzpicture}%
  }
  \caption{Diagramme de cas d'utilisation -- Sprint 4}
  \label{fig:uc_sprint4}
\end{figure}

\section{Conception}

\subsection{Cas d'utilisation raffines}

\subsubsection{UC1 -- Executer les tests de regression}
\begin{table}[H]
  \centering
  \caption{Description du cas d'utilisation UC1}
  \begin{tabularx}{\textwidth}{>{\bfseries}p{4cm} X}
    \toprule
    Cas d'utilisation & Executer les tests de regression \\
    \midrule
    Acteur principal & Equipe de developpement \\
    Preconditions & Environnement stable et donnees de test. \\
    Scenario nominal & Execution des scenarii prioritaires et collecte des resultats. \\
    Postconditions & Rapport de validation et liste d'anomalies. \\
    Exceptions & Echec de tests critiques. \\
    \bottomrule
  \end{tabularx}
\end{table}

\subsubsection{UC2 -- Deployer et stabiliser}
\begin{table}[H]
  \centering
  \caption{Description du cas d'utilisation UC2}
  \begin{tabularx}{\textwidth}{>{\bfseries}p{4cm} X}
    \toprule
    Cas d'utilisation & Deployer et stabiliser \\
    \midrule
    Acteur principal & Administrateur \\
    Preconditions & Configuration finale et verification des dependances. \\
    Scenario nominal & Deploiement, tests post-deploiement et monitoring. \\
    Postconditions & Version stable disponible pour demonstration. \\
    Exceptions & Echec de demarrage ou de connexion aux services. \\
    \bottomrule
  \end{tabularx}
\end{table}

\section{Diagrammes de sequence detailles du Sprint 4}

\subsection{Sequence detaillee -- Validation de regression}

\begin{figure}[H]
  \centering
  \begin{tikzpicture}[x=1pt, y=1pt,
    box/.style={draw, fill=gray!10, minimum width=1.8cm, minimum height=0.6cm, font=\footnotesize},
    msg/.style={draw, ->, >=Stealth, font=\scriptsize},
    ret/.style={draw, dashed, ->, >=Stealth, font=\scriptsize},
    life/.style={draw, dashed, thin}
  ]
  \foreach \x/\lbl in
    {0/Equipe, 90/Frontend, 190/API, 300/Base, 410/Logs}{
    \node[box] at (\x,0) {\lbl};
    \draw[life] (\x,-8) -- (\x,-260);
  }
  \draw[msg] (0,-30)  -- node[above]{lancer scenarios} (90,-30);
  \draw[msg] (90,-55) -- node[above]{requete API} (190,-55);
  \draw[msg] (190,-80) -- node[above]{lecture/ecriture} (300,-80);
  \draw[ret] (300,-100) -- node[above]{OK / erreurs} (190,-100);
  \draw[ret] (190,-125) -- node[above]{codes HTTP} (90,-125);
  \draw[msg] (90,-150) -- node[above]{journalisation} (410,-150);
  \draw[ret] (90,-175) -- node[above]{rapport tests} (0,-175);
  \end{tikzpicture}
  \caption{Diagramme de sequence detaille -- Validation de regression}
  \label{fig:seq_regression}
\end{figure}

\subsection{Sequence detaillee -- Deploiement}

\begin{figure}[H]
  \centering
  \begin{tikzpicture}[x=1pt, y=1pt,
    box/.style={draw, fill=gray!10, minimum width=1.8cm, minimum height=0.6cm, font=\footnotesize},
    msg/.style={draw, ->, >=Stealth, font=\scriptsize},
    ret/.style={draw, dashed, ->, >=Stealth, font=\scriptsize},
    life/.style={draw, dashed, thin}
  ]
  \foreach \x/\lbl in
    {0/Admin, 90/Serveur, 190/Backend, 300/DB, 410/Frontend}{
    \node[box] at (\x,0) {\lbl};
    \draw[life] (\x,-8) -- (\x,-260);
  }
  \draw[msg] (0,-30)  -- node[above]{deployer} (90,-30);
  \draw[msg] (90,-55) -- node[above]{demarrer API} (190,-55);
  \draw[msg] (190,-80) -- node[above]{connexion DB} (300,-80);
  \draw[ret] (300,-100) -- node[above]{OK} (190,-100);
  \draw[msg] (90,-125) -- node[above]{servir UI} (410,-125);
  \draw[ret] (410,-150) -- node[above]{OK} (90,-150);
  \draw[ret] (90,-175) -- node[above]{deploiement termine} (0,-175);
  \end{tikzpicture}
  \caption{Diagramme de sequence detaille -- Deploiement}
  \label{fig:seq_deploy}
\end{figure}

\section{Evaluation du Sprint 4}

\subsection{Scenarios de tests realises}

\begin{table}[H]
  \centering
  \caption{Table de tests fonctionnels du Sprint 4}
  \label{tab:tests_s4}
  \begin{tabular}{|p{5.2cm}|p{3.2cm}|p{5.4cm}|}
    \hline
    \textbf{Scenario} & \textbf{Resultat attendu} & \textbf{Statut} \\
    \hline
    Tests de regression complets & Tous les scenarii critiques valides & Valide \\
    \hline
    Verification roles et acces & Acces conforme aux droits & Valide \\
    \hline
    Temps de reponse API & Temps stable sur flux principaux & Valide \\
    \hline
    Deploiement local de demo & Application demarre sans erreur & Valide \\
    \hline
  \end{tabular}
\end{table}

\begin{table}[H]
  \centering
  \caption{Bilan du Sprint 4}
  \label{tab:bilan_s4}
  \begin{tabular}{|p{5cm}|c|p{6.5cm}|}
    \hline
    \textbf{Fonctionnalite} & \textbf{Statut} & \textbf{Observations} \\
    \hline
    Integration globale & Fait & Flux complets verifies \\
    \hline
    Correction anomalies & Fait & Bugs critiques resolus \\
    \hline
    Validation securite & Fait & Roles et droits verifies \\
    \hline
    Stabilisation deploiement & Fait & Demo stable \\
    \hline
    Documentation finale & Fait & Technique + utilisateur \\
    \hline
  \end{tabular}
\end{table}

\section*{Conclusion}
\addcontentsline{toc}{section}{Conclusion}

Le Sprint 4 a permis la validation finale du systeme, avec une integration
stable, des tests complets, un deploiement pret pour la demonstration et une
documentation finalisee.

\clearpage

% ============================================================
\end{document}
% ============================================================
